% !TeX root = ../../Book5.tex
% 纲要标题：## 附录 E　微分模与微分 Galois 理论
% 纲要位置：工作记录/计划纲要.md，第 5058 行。
% > 所需基础是第三卷附录 E.1 的微分域语言、模与张量，以及本卷第21章中所需的线性代数群语言；不要求完整非线性消去、半单 Lie 代数分类或导出范畴。

\chapter{微分模与微分 Galois 理论}
\label{b5:app:E}
\BookChapterNumber{b5:app:E}

\begin{chapterabstract}
本附录从带导子的域与微分模建立 Picard--Vessiot 扩张，研究无新常数条件、微分 Galois 群及基本对应，并计算指数、对数和根式型扩张。
\end{chapterabstract}

\begin{learninggoals}
\begin{itemize}
\item 能在微分模、线性系统及微分算子商模之间转换，并正确处理对偶的符号。
\item 检验基本解矩阵和无新常数条件，验证具体 PV 扩张并求出其全部自同构。
\item 通过解的代数关系区分乘法群、加法群、有限群及其子群。
\item 使用单位连通分支的可解性准则，计算张量与行列式对群的限制。
\end{itemize}
\end{learninggoals}

在微分域 \((\mathbb C(x),\mathrm d/\mathrm dx)\) 上添入 \(u'=u\) 的一个非零解，并要求扩张不增加常数。固定 \(\mathbb C(x)\) 且与求导交换的自同构 \(\sigma\) 使 \((\sigma u)'=\sigma(u')=\sigma u\)，故 \((\sigma u/u)'=0\)。无新常数的条件使这个比值属于 \(\mathbb C\)，又因 \(\sigma u\ne0\)，只能有 \(\sigma u=cu\)、\(c\in\mathbb C^\times\)。多项式根的 Galois 群由有限置换描述，这里却自然出现一个代数群。线性微分方程的对称性为何有这种形式，保持全部代数关系又如何限制常数矩阵？本附录由具体扩张进入 Picard--Vessiot 理论与微分 Galois 对应。

\ProseParagraphBreak
需要特征零域扩张、左模、张量积和矩阵代数；存在唯一性的证明还用到 Noether 规范化、整扩张与弱 Nullstellensatz。Galois 对应及重构使用线性代数群、Zariski 闭子群和有理表示。第三卷附录 E.1 的导子语言在此重新给出。

\ProseParagraphBreak
Singer 的综述收入 2009 年的《Algebraic Theory of Differential Equations》\cite{Singer2009DifferentialGalois}，下文所引定理号及页码依作者公开稿。较详细的证明参见 van der Put 与 Singer 的专著\cite{VanDerPutSinger2003GaloisTheory}，本附录对它的编号按作者 2002 年修订公开稿核对。

% !TeX root = ../../Book5.tex
% 纲要标题：### E.1 微分模与线性方程
% 纲要位置：工作记录/计划纲要.md，第 5062 行。
% 纲要细目（保留原定写作范围）：
%
% * 微分域上的有限维微分模、Leibniz 条件及选基后的系统 \(\delta Y=AY\)
% * 微分算子环 \(K[\partial;\delta]\) 的模、张量、对偶与水平截面；联络算子（connection operator）不要求平方为零
% * 与第二卷 Ore 扩张及本卷附录 D 的 D-模比较；二阶方程模的对偶恢复伴随系统，迹的换基与行列式模相联系
% #### E.1.1 微分模、Leibniz 条件与线性系统
% #### E.1.2 微分算子环、张量与水平截面
% #### E.1.3 Ore 扩张与 D-模的低维比较
%
\section{微分模与线性方程}
\label{b5:sec:E01}

线性微分方程的未知量可以写成函数，也可以组织成向量空间上的附加算子。后者允许我们谈论子模、张量积和对偶；换基时变化的是方程矩阵，方程所描述的微分模则不变。本节固定一个导子，不讨论若干偏导算子的相容性。

\subsection{微分模、Leibniz 条件与线性系统}
\label{b5:subsec:E01:01}

\begin{definition}\label{b5:def:E-differential-module}
	设 \(K\) 为特征零域，\(\delta:K\to K\) 为加法映射，满足 \(\delta(ab)=\delta(a)b+a\cdot\delta(b)\)。二元组 \((K,\delta)\) 称为\term[weifen-yu]{微分域}[differential field]，其\term[changshu-yu]{常数域}[field of constants]为 \(C_K=\ker\delta\)。若有限维 \(K\) 向量空间 \(M\) 上的加法算子 \(\nabla\) 满足
	\[
		\nabla(am)=\delta(a)m+a\nabla(m)\qquad(a\in K,m\in M),
	\]
	则 \((M,\nabla)\) 称为 \((K,\delta)\) 上的\term[weifen-mo]{微分模}[differential module]。对两个微分模 \((M,\nabla_M)\)、\((N,\nabla_N)\)，若 \(K\) 线性映射 \(T:M\to N\) 满足
	\[
		T\circ\nabla_M=\nabla_N\circ T,
	\]
	则称 \(T\) 为微分模同态。
\end{definition}
\symidx{\((K,\delta),C_K\)：微分域与其常数域}
\symidx{\(\nabla\)：微分模上满足 Leibniz 法则的联络算子}

常数域确实是域：由 \(\delta(1)=0\) 和 \(\delta(a^{-1})=-a^{-2}\cdot\delta(a)\)，非零常数的逆仍为常数。第三卷命题 E.1.2 证明了整环上的导子向分式域的唯一延拓；这里所需的公式就是
\[
	\delta(a/b)=\frac{\delta(a)b-a\cdot\delta(b)}{b^2}.
\]
\(\nabla\) 一般只对常数域线性。它是联络方向上的算子，不是复形的微分，定义不要求 \(\nabla^2=0\)。

\ProseParagraphBreak
把一组基写成行 \(e=(e_1,\ldots,e_n)\)，约定 \(\nabla(e_j)=-\sum_i A_{ij}e_i\)。对坐标列 \(v\) 有
\begin{equation}\label{b5:eq:E-horizontal-system}
	\nabla(ev)=e\bigl(\delta(v)-Av\bigr).
\end{equation}
所以 \(\nabla(ev)=0\) 恰好是线性系统 \(\delta(v)=Av\)。这里的负号使水平向量与通常写法的解相符。 若新基 \(\widetilde e=eP\)，其中 \(P\in\operatorname{GL}_n(K)\)，则直接应用 Leibniz 法则得到
\begin{equation}\label{b5:eq:E-gauge-change}
	\widetilde A=P^{-1}AP-P^{-1}\cdot\delta(P).
\end{equation}
这称为\term[guifan-bianhuan]{规范变换}[gauge transformation]；它不是通常的矩阵相似，因为基向量的系数也要被求导。

对该式取迹，利用行列式的求导公式可得
\[
	\operatorname{tr}\widetilde A =\operatorname{tr}A-\operatorname{tr}\bigl(P^{-1}\cdot\delta(P)\bigr) =\operatorname{tr}A-\frac{\delta(\det P)}{\det P}.
\]
迹虽然随基改变，却只相差一个对数导数。若 \(\omega=e_1\wedge\cdots\wedge e_n\)，张量上的 Leibniz 法则在最高外幂中给出 \(\nabla\omega=-\operatorname{tr}(A)\omega\)：非对角项会产生重复的基向量，只有对角项留下。因此秩一微分模 \(\bigwedge^nM\) 的水平方程为 \(\delta(f)=\operatorname{tr}(A)f\)。它有基域内的水平基，当且仅当 \(\operatorname{tr}A=\delta(f)/f\) 对某个 \(f\in K^\times\) 成立；换基公式说明这个条件与基的选择无关。\cref{b5:prop:E-wronskian}将把同一个秩一方程写成基本解矩阵的 Wronskian 方程。

\ProseParagraphBreak
例如 \(K=\mathbb C(x)\)、\(\delta=\mathrm{d}/\mathrm{d}x\) 上的二阶方程 \(y''+ay'+by=0\) 等价于
\[
	\binom{y}{y'}'=\begin{pmatrix}0&1\\-b&-a\end{pmatrix}\binom{y}{y'}.
\]
第一行保证第二坐标是第一坐标的导数，第二行再恢复原方程，因而两个方向都成立。一般系统化成单个标量方程，要使用微分模的循环向量定理：当 \(\delta\ne0\) 时，每个秩 \(n\) 微分模都有向量 \(v\)，使 \(v,\nabla v,\ldots,\nabla^{n-1}v\) 为 \(K\) 基。因而它可写成一个循环 \(K[\partial;\delta]\) 模。把这一结果用于系统所对应模的对偶，便恢复标量方程的解系统；见\cite[Theorem 1.2.8 and Corollary 1.2.9, pp. 6--7]{Singer2009DifferentialGalois}及\cite[Section 2.1]{VanDerPutSinger2003GaloisTheory}。这里的 \(\nabla\) 满足 Leibniz 法则而不是 \(K\) 线性算子，不能直接用普通线性算子的有理标准形替代这个定理。

\subsection{微分算子环、张量与水平截面}
\label{b5:subsec:E01:02}

形式多项式 \(\sum_{j=0}^m a_j\partial^j\) 在关系
\[
	\partial a=a\partial+\delta(a)
\]
下组成环 \(K[\partial;\delta]\)，系数一律写在左边。第二卷定理 20.5.2 证明了 Ore 扩张\footnote{奥尔（Øystein Ore，1899-1968），挪威数学家，研究环论、Galois 理论和图论；其非交换多项式环以他命名。}乘法的存在与结合性；这里取底环 \(K\)、自同态 \(\sigma=\operatorname{id}_K\)，其 \(\sigma\) 导子条件恰为通常的 Leibniz 法则。微分模对应于在 \(K\) 上有限维的左 \(K[\partial;\delta]\) 模：令 \(\partial m=\nabla(m)\)，上面的环关系正好成为定义中的 Leibniz 法则。
\symidx{\(K[\partial;\delta]\)：微分域的 Ore 微分算子环}

\begin{proposition}\label{b5:prop:E-tensor-horizontal}
	设 \((M,\nabla_M)\)、\((N,\nabla_N)\) 为特征零微分域 \((K,\delta)\) 上的有限维微分模，\(C_K=\ker\delta\)，\(M^*=\operatorname{Hom}_K(M,K)\)。下列公式分别在 \(M\otimes_KN\) 和 \(M^*\) 上定义微分模结构：
	\[
		\begin{aligned}
			\nabla(m\otimes n)&=\nabla_M(m)\otimes n+m\otimes\nabla_N(n),\\
			(\nabla\varphi)(m)&=\delta\bigl(\varphi(m)\bigr)-\varphi\bigl(\nabla_M(m)\bigr) \qquad(\varphi\in M^*).
		\end{aligned}
	\]
	若 \(H(M)=\ker\nabla_M\)，则 \(H(M)\) 是 \(C_K\) 向量空间，且 \(\dim_{C_K}H(M)\le\dim_KM\)。
\end{proposition}
\symidx{\(H(M)=\ker\nabla_M\)：微分模的水平截面空间}
\begin{proof}
	将 \((am)\otimes n\) 代入第一式得到的额外项是 \(\delta(a)m\otimes n\)，将 \(m\otimes(an)\) 代入也得到同一项，故公式尊重张量平衡关系。再把 \(a\) 提到任一因子外即可验证 Leibniz 法则。对偶公式中，将 \(am\) 代入时两个 \(\delta(a)\varphi(m)\) 项抵消，故 \(\nabla\varphi\) 仍为 \(K\) 线性泛函；对 \(a\varphi\) 代入则给出相应 Leibniz 法则。

	最后证明任意常数线性无关的水平向量也在 \(K\) 上线性无关。若不然，取项数最少的非平凡关系，归一化为 \(m_r=\sum_{i<r}a_i m_i\)。其前 \(r-1\) 项在 \(K\) 上无关；对关系施加 \(\nabla\)，得到 \(\sum_{i<r}\delta(a_i)m_i=0\)，从而所有 \(a_i\) 都是常数，与原来的常数线性无关矛盾。维数上界随即成立。
\end{proof}

\begin{proposition}\label{b5:prop:E-dual-system}
	设 \((K,\delta)\) 为特征零微分域，\(M\) 为有限维微分 \(K\) 模，在基 \(e=(e_1,\ldots,e_n)\) 下满足 \(\nabla(e)=-eA\)，其中 \(A\in M_n(K)\)。在对偶基下，\(M^*\) 的水平坐标方程为 \(\delta(w)=-A^{\mathsf T}w\)。在任意微分扩域中，若 \(\delta(v)=Av\) 且 \(\delta(w)=-A^{\mathsf T}w\)，则配对 \(w^{\mathsf T}v\) 为常数。
\end{proposition}
\begin{proof}
	对偶联络公式给出 \((\nabla e_i^*)(e_j)=-e_i^*(\nabla e_j)=A_{ij}\)，所以 \(\nabla e_i^*=\sum_jA_{ij}e_j^*\)。若仍按本节约定写 \(\nabla(e^*)=-e^*A_{\mathrm{dual}}\)，则 \(A_{\mathrm{dual}}=-A^{\mathsf T}\)；相应水平系统为 \(\delta(w)=-A^{\mathsf T}w\)。对配对求导，两项恰好抵消：
	\[
		\delta(w^{\mathsf T}v)=-w^{\mathsf T}Av+w^{\mathsf T}Av=0.
	\]
\end{proof}

\(H(M)\) 中的元素称为\term[shuiping-jiemian]{水平截面}[horizontal section]。在微分扩域 \(L/K\) 中，延拓公式为 \(\nabla(\ell\otimes m)=\delta(\ell)\otimes m+\ell\otimes\nabla(m)\)。 延拓后可能出现足够多的水平截面。例如 \(\nabla(e)=-e\) 在 \(\mathbb C(x)e\) 中没有非零水平截面，而加入 \(t\) 满足 \(t'=t\) 后，\(te\) 就是水平基向量。

\subsection{Ore 扩张与 D-模的低维比较}
\label{b5:subsec:E01:03}

在 \(K=\mathbb C(x)\) 中，\(D=K[\partial;\mathrm{d}/\mathrm{d}x]\) 含有多项式系数的 Weyl 代数 \(A_1(\mathbb C)\)。但 \(D\) 模忘记了有限点上的部分信息：在 \(A_1\) 模 \(A_1/A_1x\) 中，非零向量 \(\overline1\) 被 \(x\) 零化，扩系数可写为 \(D\otimes_{A_1}(A_1/A_1x)\simeq D/Dx\)。在 \(D\) 中 \(x\) 可逆，故生成元满足 \(\overline1=x^{-1}\cdot x\overline1=0\)，整个 \(D\) 模也变为零。这样的点支撑模不可能由一个非零的有限维微分域模反映。

\ProseParagraphBreak
还应区别两种“解”。非零标量算子 \(P\in D\) 给出左模 \(Q=D/DP\)，它在 \(K\) 上有限维：若 \(P\) 的阶为 \(r\)，利用可逆首系数逐次减去 \(P\) 的左倍数，便得到阶小于 \(r\) 的唯一余式；唯一性来自非零乘积的阶相加。因此 \(1,\partial,\ldots,\partial^{r-1}\) 的类为基（\(r=0\) 时商模为零）。对任意微分扩域 \(L/K\)，映射
\[
	\operatorname{Hom}_D(Q,L)\longrightarrow\bigl\{y\in L:P(y)=0\bigr\}, \qquad f\longmapsto f(\overline1)
\]
是双射：从解 \(y\) 反向定义 \(f(R+DP)=R(y)\)，关系 \(P(y)=0\) 保证良定义。这是反变的解空间。与之相对，\eqref{b5:eq:E-horizontal-system}用的是指定微分模的水平截面；前者等于对偶微分模的水平截面。若 \(P=\partial-a\)，则 \(\partial\overline1=a\overline1\)，该商模本身的水平坐标满足 \(v'=-av\)，而其同态解满足 \(y'=ay\)。这里的符号差别来自对偶。取 \(a=\lambda/x\)，则 \(Q\) 自身的水平坐标为 \(cx^{-\lambda}\)，同态解为 \(cx^\lambda\)，在正向解析延拓下分别乘以 \(e^{-2\pi i\lambda}\) 与 \(e^{2\pi i\lambda}\)。附录 D 的\subsecref{b5:subsec:D02:04}用 \(\lambda=1/3\) 比较了这两个局部系统。按本节 \(\nabla(e)=-eA\) 的约定，取 \(A=\lambda/x\) 得到的系统模是 \(Q^*\)，所以其水平截面才直接满足原标量方程。

\ProseParagraphBreak
对开头的二阶算子 \(P=\partial^2+a\partial+b\)，取商模基 \(e_0=\overline1,e_1=\overline\partial\)。左乘作用给出 \(\nabla e_0=e_1\)、\(\nabla e_1=-be_0-ae_1\)，故按 \(\nabla(e)=-eA_Q\) 的约定，
\[
	A_Q=\begin{pmatrix}0&b\\-1&a\end{pmatrix},
	\qquad -A_Q^{\mathsf T}=\begin{pmatrix}0&1\\-b&-a\end{pmatrix}.
\]
\cref{b5:prop:E-dual-system}因此恢复了原方程的一阶系统。具体地，\(D\) 线性同态 \(f:Q\to L\) 在这组基上的值为 \((f(e_0),f(e_1))=(y,y')\)；这些值是对偶模的水平坐标，而非商模自身的水平坐标。

% !TeX root = ../../Book5.tex
% 纲要标题：### E.2 Picard–Vessiot 扩张与对应
% 纲要位置：工作记录/计划纲要.md，第 5074 行。
% 纲要细目（保留原定写作范围）：
%
% * 在特征零、常数域代数闭的标准情形下介绍 PV 环、基本解矩阵、微分单性与无新常数条件
% * 保持基域且与导子交换的自同构构成常数域上的线性代数群；证明存在唯一性及群闭性
% * 中间微分域与 Zariski 闭子群的对应；闭性条件不能删去，给出证明概要。指数比值与对数差值展示主齐性空间坐标及余乘法，并区分普通中间域与中间微分域
% #### E.2.1 PV 环、基本解矩阵与无新常数条件
% #### E.2.2 微分自同构群与存在唯一性
% #### E.2.3 中间微分域与 Zariski 闭子群
%
\section{Picard–Vessiot 扩张与对应}
\label{b5:sec:E02}

一个方程可能没有足够多的有理函数解。为研究全部线性解，应当加入基本解矩阵，同时控制新增的常数。这里始终固定特征零微分域 \((K,\delta)\)，并要求常数域 \(C=\ker\delta\) 代数闭。这个假设保证下述存在唯一性具有最简洁的形式。

\subsection{PV 环、基本解矩阵与无新常数条件}
\label{b5:subsec:E02:01}

\begin{definition}\label{b5:def:E-pv-ring}
设 \((K,\delta)\) 为特征零微分域，常数域 \(C=\ker\delta\) 代数闭，\(n\ge1\)，\(A\in M_n(K)\)。设交换微分 \(K\) 代数 \(R\) 是整环。若存在 \(Z\in\operatorname{GL}_n(R)\) 满足 \(\delta(Z)=AZ\)，有 \(R=K\bigl[Z_{ij},(\det Z)^{-1}\bigr]\)，且 \(R\) 的微分理想只有 \(0,R\)，则称 \(R\) 为系统 \(\delta(Y)=AY\) 的\term[Picard-Vessiot-huan]{Picard–Vessiot 环}[Picard--Vessiot ring]，简称 PV 环。这里微分理想指满足 \(\delta(I)\subset I\) 的理想；\(Z\) 称为\term[jiben-jie-juzhen]{基本解矩阵}[fundamental solution matrix]。分式域 \(L=\operatorname{Frac}R\) 称为相应的\term[Picard-Vessiot-kuozhang]{Picard–Vessiot 扩张}[Picard--Vessiot extension]。
\end{definition}

微分单性不等于环是域：稍后会出现 \(K[t,t^{-1}]\) 这样的 PV 环。其作用是排除与方程无关的微分商；而在上述常数假设下，PV 分式域满足 \(C_L=C\)，这称为无新常数。反过来，由基本解矩阵的系数生成、且无新常数的微分域扩张，也给出这样的 PV 环。下一小节证明这些性质。

\ProseParagraphBreak
若只要求存在解矩阵，就可以把 \(Z\) 乘以任意新增的超越常数，因而无法确定扩张。无新常数限制的是全部分式域，不能只检查所选解的几个系数。基本解矩阵也不是逐列任取非零解：必须要求行列式非零。

\subsection{微分自同构群与存在唯一性}
\label{b5:subsec:E02:02}

取微分理想的商，可以消去独立解变量之间不应保留的常数自由度。例如在 \(K=\mathbb C(x)\) 上，先令 \(U,V\) 为代数独立元，规定 \(U'=U,V'=2V\)。在
\(S=K[U^{\pm1},V^{\pm1}]\) 的分式域中有
\[
 \left(\frac V{U^2}\right)'=0,
 \qquad \frac V{U^2}\notin\mathbb C.
\]
所以 \(\operatorname{diag}(U,V)\) 虽是可逆解矩阵，这个扩张却有新常数。若取 \(a\in\mathbb C^\times\)，关系 \(V-aU^2\) 的导数为 \(2(V-aU^2)\)，因此 \((V-aU^2)\) 是微分理想。商环
\(S/(V-aU^2)\simeq K[U,U^{-1}]\) 将上述常数固定为 \(a\)；\cref{b5:exm:E-exponential}会证明它无新常数且为 PV 环。改变 \(a\) 只须相应地伸缩 \(V\)，所得微分 \(K\) 代数同构。一般存在性证明同样要对通用解变量施加微分关系，直到商环微分单；下面的引理保证，此时不会留下新常数。

\begin{lemma}\label{b5:lem:E-simple-constants}
设 \((K,\delta)\) 为特征零微分域，常数域 \(C=\ker\delta\) 代数闭。若非零交换微分 \(K\) 代数 \(S\) 作为 \(K\) 代数有限生成，且其微分理想只有 \(0,S\)，则 \(S\) 是整环，且 \(\operatorname{Frac}S\) 的常数域为 \(C\)。
\end{lemma}
\begin{proof}
若 \(a\in S\) 不幂零，则
\(I_a=\{b:a^mb=0\;\text{对某个}\;m\ge1\}\)
是一个真微分理想：对 \(a^mb=0\) 求导再乘 \(a\)，得到 \(a^{m+1}b'=0\)。故 \(I_a=0\)，\(a\) 不是零因子。若非零 \(a\) 幂零，取最小 \(m\) 使 \(a^m=0\)，则
\(ma^{m-1}a'=0\) 且 \(ma^{m-1}\ne0\)，所以 \(a'\) 是零因子，按刚才的结论必幂零。于是幂零元构成微分理想，只能为零。这证明 \(S\) 为整环。

若 \(c\in\operatorname{Frac}S\) 满足 \(c'=0\)，非零理想
\(\{b\in S:bc\in S\}\) 对导子稳定，故为 \(S\)，从而 \(c\in S\)。若 \(c\notin C\)，则对每个 \(a\in C\)，\((c-a)S\) 是非零微分理想，所以 \(c-a\) 都可逆。

这迫使 \(c\) 在 \(K\) 上代数。确实，若 \(c\) 超越，在 \(K(c)\) 上对 \(S\otimes_{K[c]}K(c)\) 应用第三卷第5.4节的 Noether 正规化，得到代数独立元 \(u_1,\ldots,u_d\)，使该代数在 \(K(c)[u_1,\ldots,u_d]\) 上有限。先清去 \(u_i\) 表达式的分母，再清去 \(S\) 的有限个代数生成元所满足的首一整性方程中系数的分母。这些分母只有有限多个，其乘积给出非零 \(f\in K[T]\)，于是 \(u_i\in S[f(c)^{-1}]\)，且 \(S[f(c)^{-1}]\) 整于多项式环
\[
B=K[c,f(c)^{-1}][u_1,\ldots,u_d].
\]
这里 \(u_i\) 在 \(K(c)\) 上代数独立。常数域 \(C\) 无限，而 \(f\) 只有有限多个根，故可取 \(a\in C\) 使 \(f(a)\ne0\)。于是
\[
B/(c-a)B\simeq K[u_1,\ldots,u_d],
\]
右侧为非零整环，所以 \((c-a)B\) 是真素理想。第三卷第5.2节的 Lying-over 定理用于这个整扩张，给出 \(S[f(c)^{-1}]\) 中收缩为它的素理想，该素理想包含 \(c-a\)。可是 \(c-a\) 已在 \(S\) 中可逆，在局部化中仍为单位，不可能属于真素理想，矛盾。最后，\(c\) 的首一最小多项式求导后，最高项消失；最小性说明其全部系数的导数为零。该多项式因此在 \(C[T]\) 中，\(C\) 代数闭又推出 \(c\in C\)。
\end{proof}

\begin{theorem}[Picard–Vessiot 存在唯一性]\label{b5:thm:E-pv-existence}
设 \((K,\delta)\) 为特征零微分域，常数域 \(C=\ker\delta\) 代数闭，\(n\ge1\)，\(A\in M_n(K)\)。则系统 \(\delta(Y)=AY\) 存在 PV 环，任意两个这样的环在微分 \(K\) 代数意义下同构；其分式域的常数域均为 \(C\)。反之，若 \(L/K\) 是微分域扩张，\(Z\in\operatorname{GL}_n(L)\) 满足 \(\delta(Z)=AZ\)，\(L=K(Z_{ij})\) 且 \(C_L=C\)，则 \(K\bigl[Z_{ij},(\det Z)^{-1}\bigr]\) 是该系统的 PV 环。
\end{theorem}

\begin{proof}
在 \(S_0=K[X_{ij},(\det X)^{-1}]\) 上按 \(X'=AX\) 延伸导子，取一个极大真微分理想 \(I\)。商环 \(R=S_0/I\) 微分单，且仍有可逆基本解矩阵。由\cref{b5:lem:E-simple-constants}，它是整环，分式域没有新常数，所以是所需 PV 环。

若 \(R_1,R_2\) 是两个 PV 环，在 \(R_1\otimes_KR_2\) 中取极大真微分理想并作商，得到微分单环 \(T\)。两个 \(R_i\to T\) 都单射，且由同一引理 \(T\) 的常数为 \(C\)。两基本解矩阵在 \(T\) 中的比值导数为零，故相差一个 \(\operatorname{GL}_n(C)\) 矩阵。它们的系数及逆行列式因而生成同一个子环，于是两个嵌入的像相等，给出所需同构。

为证反向陈述，需要一个常数变量引理：若微分域 \(F\) 的常数为 \(C\)，且变量 \(U_{ij}\) 的导数全为零，则
\(F[U_{ij},(\det U)^{-1}]\) 中的每个微分理想都由它与 \(C[U_{ij},(\det U)^{-1}]\) 的交生成。记 \(J\) 为所考察的微分理想，并令
\[
 B=C[U_{ij},(\det U)^{-1}],\qquad
 J_0=(J\cap B)\,F[U_{ij},(\det U)^{-1}],
\]
取 \(B\) 的一组 \(C\) 基 \((e_\alpha)\)，把 \(p\in J\) 唯一写成有限项 \(p=\sum a_\alpha e_\alpha\)，按非零系数个数归纳证明 \(p\in J_0\)。零项时结论显然；非零时先除以一个非零系数，使其中一个系数为 \(1\)。若全部系数都在 \(C\) 中，则归一化后的 \(p\in J\cap B\)，结论成立。否则取另一个系数 \(a\notin C\)。因为 \(J\) 对导子稳定，
\[
 p'\in J,\qquad
 \left(\frac pa\right)'=\frac{p'}a-\frac{a'}{a^2}p\in J.
\]
在 \(p'\) 中，原来系数为 \(1\) 的项消失；在 \((p/a)'\) 中，原来系数为 \(a\) 的项消失。两个导数的非零系数个数都严格少于 \(p\)，故归纳假设使它们都属于 \(J_0\)。于是
\[
 \frac{a'}{a^2}p=\frac{p'}a-\left(\frac pa\right)'\in J_0.
\]
由 \(C=\ker\delta\) 和 \(a\notin C\) 得 \(a'\ne0\)，所以 \(a'/a^2\) 在域 \(F\) 中可逆，推出 \(p\in J_0\)。乘回最初用于归一化的系数，便得到原理想元也属于 \(J_0\)。

现设 \(L=K(Z_{ij})\)、\(Z'=AZ\)、\(C_L=C\)。仍取 \(S_0\) 的极大真微分理想 \(I\)，在
\(L\otimes_KS_0\) 中改变量为 \(X=ZU\)。因为 \(Z\) 可逆，这给出变量环的可逆变换，且
\[
 U'=(Z^{-1}X)'=-Z^{-1}AZZ^{-1}X+Z^{-1}AX=0.
\]
因此 \(U\) 的各条目成为常数变量，刚才的引理把 \(I\) 的延伸化成由常数多项式理想 \(J\) 生成的理想。它是真理想，因为 \(L\otimes_K(S_0/I)\ne0\)。由弱 Nullstellensatz，\(J\) 有一个可逆常数矩阵零点 \(c\)。代入 \(X=Zc\) 得到微分同态 \(S_0\to L\)，其核包含 \(I\)；极大微分性使二者相等。其像正好是
\(K[Z_{ij},(\det Z)^{-1}]\)，故此环微分单，证明反向陈述。完整证明的相同代数构造见\cite[Section 1.3, Lemmas 1.17 and 1.23, Propositions 1.20 and 1.22]{VanDerPutSinger2003GaloisTheory}。
\end{proof}

同构一般不唯一，其不唯一性正是下一定义要研究的对象。概念及存在唯一性的概述见\cite[Definition 1.3.2 and the following discussion, pp. 10--11]{Singer2009DifferentialGalois}。

\begin{definition}\label{b5:def:E-galois-group}
设 \((K,\delta)\) 为特征零微分域，常数域 \(C=\ker\delta\) 代数闭，\(L/K\) 为 PV 扩张。所有固定 \(K\) 且满足 \(\sigma\delta=\delta\sigma\) 的 \(L\) 的域自同构组成群
\[
 G=\operatorname{Gal}_{\delta}(L/K).
\]
这个群称为该扩张的\term[weifen-Galois-qun]{微分 Galois 群}[differential Galois group]。
\end{definition}
\symidx{\(\operatorname{Gal}_\delta(L/K)\)：Picard--Vessiot 扩张的微分 Galois 群}

\begin{proposition}\label{b5:prop:E-group-matrix}
设 \((K,\delta)\) 为特征零微分域，常数域 \(C=\ker\delta\) 代数闭，\(n\ge1\)，\(A\in M_n(K)\)。设 \(L/K\) 为系统 \(\delta(Y)=AY\) 的 PV 扩张，\(Z\in\operatorname{GL}_n(L)\) 为基本解矩阵。每个 \(\sigma\in\operatorname{Gal}_\delta(L/K)\) 唯一写成 \(\sigma(Z)=ZC_\sigma\)，其中 \(C_\sigma\in\operatorname{GL}_n(C)\)；对应 \(\sigma\mapsto C_\sigma\) 是忠实群表示。换用另一基本解矩阵只使此表示共轭。
\end{proposition}
\begin{proof}
利用 \(\delta(Z^{-1})=-Z^{-1}\cdot\delta(Z)Z^{-1}\)，直接求导得到 \(\delta\bigl(Z^{-1}\cdot\sigma(Z)\bigr)=0\)。无新常数保证该矩阵在 \(C\) 上。因为自同构固定常数，\(\sigma\tau(Z)=ZC_\sigma C_\tau\)，所以这是群同态。\(Z\) 的系数生成 \(L\)，故核平凡。若 \(Z'\) 是另一基本解矩阵，同一求导计算给出 \(Z'=ZB\)、\(B\in\operatorname{GL}_n(C)\)，于是新矩阵为 \(B^{-1}C_\sigma B\)。
\end{proof}

\begin{corollary}\label{b5:cor:E-group-closed}
设 \((K,\delta)\) 为特征零微分域，常数域 \(C=\ker\delta\) 代数闭，\(n\ge1\)，\(A\in M_n(K)\)，\(L/K\) 为系统 \(\delta(Y)=AY\) 的 PV 扩张，\(Z\in\operatorname{GL}_n(L)\) 为基本解矩阵。由 \(\sigma(Z)=ZC_\sigma\) 给出的微分 Galois 群矩阵像，是 \(\operatorname{GL}_n(C)\) 中的 Zariski 闭子群。
\end{corollary}
\begin{proof}
将 PV 环写成 \(R=K[X_{ij},(\det X)^{-1}]/I\)，其中 \(I\) 是极大真微分理想。取有限个理想生成元 \(f_1,\ldots,f_r\)。常数矩阵 \(c\) 给出自同构，当且仅当各个 \(f_j(Zc)=0\)：此条件让 \(X\mapsto Zc\) 通过商环；所得微分同态的像含 \(Z=(Zc)c^{-1}\)，所以满射，微分单性又保证单射。

对每个 \(j\)，把 \(f_j(Zc)\) 展开为 \(R\) 中有限多个 \(C\) 线性无关元素的组合，其系数是 \(c_{ij}\) 与 \((\det c)^{-1}\) 的多项式。它为零，当且仅当这些有限个系数全部为零。这是一组 \(\operatorname{GL}_n(C)\) 上的正则方程，因而给出闭性；其含义正是保持解矩阵的全部代数关系。亦见\cite[Proposition 1.3.6, p. 12]{Singer2009DifferentialGalois}。
\end{proof}

下一节会直接检查具体关系及全部允许的变换。

\subsection{中间微分域与 Zariski 闭子群}
\label{b5:subsec:E02:03}

\begin{theorem}[微分 Galois 对应]\label{b5:thm:E-galois-correspondence}
设 \((K,\delta)\) 为特征零微分域，常数域 \(C=\ker\delta\) 代数闭，\(L/K\) 为 PV 扩张，\(G=\operatorname{Gal}_\delta(L/K)\) 为相应的 \(C\) 上线性代数群。则
\[
 H\longmapsto L^H,\qquad
 F\longmapsto\operatorname{Gal}_\delta(L/F)
\]
给出 \(G\) 的 Zariski 闭子群与中间微分域 \(K\subset F\subset L\) 之间逆序的双射。其中 \(L^H\) 是被 \(H\) 全部元素固定的子域。对应的 \(F/K\) 为 PV 扩张当且仅当 \(H\) 正规，此时 \(\operatorname{Gal}_\delta(F/K)\cong G/H\)。
\end{theorem}

\begin{proofsketch}
记 \(R\) 为 \(L/K\) 的一个 PV 环，\(Z\in\operatorname{GL}_n(R)\) 为相应 \(K\) 系数系统 \(\delta(Z)=AZ\) 的基本解矩阵。在 \(R\otimes_KR\) 中，两基本解矩阵
\(Z_1=Z\otimes1\)、\(Z_2=1\otimes Z\) 的比较矩阵
\(U=Z_1^{-1}Z_2\) 满足 \(U'=0\)。证明\cref{b5:thm:E-pv-existence}时的常数变量引理，先在 \(L\otimes_KR\) 中把全部关系化为常数关系，再收缩回 \(R\otimes_KR\)，得到
\[
 R\otimes_KR\simeq R\otimes_C C[G],\qquad Z_2=Z_1U.
\]
这些常数关系的零点 \(c\) 给出 \(Z\mapsto Zc\) 的微分自同构，反向每个微分自同构也给出一个零点。这里是张量积代数 \(R\otimes_KR\)，不是由各次张量幂组成的张量代数。特征零下，有限生成整 \(K\) 代数 \(R\) 几何约简，所以 \(R\otimes_KR\) 约简。若将常数关系的商代数先记为 \(B\)，上式给出 \(R\otimes_C B\simeq R\otimes_KR\)。映射 \(B\to R\otimes_C B\) 单射，因此 \(B\) 也没有非零幂零元，即常数关系的理想是根理想。由 Nullstellensatz，这些关系正好定义上述自同构组成的闭群 \(G\)，故 \(B=C[G]\)。这个同构表示 \(\operatorname{Spec}R\) 是一个 \(G\) 主齐性空间：选定一组解后，另一组解由唯一的群元素联系。

例如对下一节的指数方程 \(\delta t=t\)，取 \(K=\mathbb C(x)\)、\(\delta x=1\) 及 PV 环 \(R=K[t,t^{-1}]\)，其中 \(t\) 在 \(K\) 上超越。在 \(R\otimes_KR\) 中令 \(t_1=t\otimes1\)、\(t_2=1\otimes t\)，则
\[
u=t_2/t_1,\qquad \delta u=0,\qquad
R\otimes_KR=K[t_1^{\pm1},t_2^{\pm1}]
\simeq R[u,u^{-1}].
\]
右端 \(R\) 对应第一份解变量 \(t_1\)，同构的逆把 \(t_2\) 写成 \(t_1u\)。群为 \(\mathbb G_m\)，其坐标环是 \(C[u,u^{-1}]\)：两份解的比值成为常数，恰好记录把第一份解伸缩成第二份解的群元素。一般比较矩阵 \(U\) 是这个比值的矩阵版本。

对\cref{b5:exm:E-logarithm}中的 \(R=K[u]\)、\(u'=1/x\)，则令 \(c=u_2-u_1\)，得到 \(c'=0\) 及
\[
 R\otimes_KR=K[u_1,u_2]=R[c]\simeq R\otimes_{\mathbb C}\mathbb C[c].
\]
这些比较坐标还记录群乘法。在三份解的张量积中，
\[
 c_{13}=c_{12}+c_{23},\qquad
 t_3/t_1=(t_2/t_1)(t_3/t_2).
\]
因此加法坐标的余乘法为 \(\Delta(c)=c\otimes1+1\otimes c\)，乘法坐标的余乘法为 \(\Delta(u)=u\otimes u\)。同一计算给出余单位 \(\varepsilon(c)=0,\varepsilon(u)=1\) 和逆映射 \(S(c)=-c,S(u)=u^{-1}\)。两份基本解的比较及三份比较的复合，恢复的是整个坐标 Hopf 代数，而不只是群的点集。

若 \(f\in L\) 被全部 \(G\) 元素固定，群常数点的 Zariski 稠密性使上述比较恒等式给出
\(f\otimes1=1\otimes f\) 在 \(L\otimes_KL\) 中成立。域扩张的忠实平坦下降于是推出 \(f\in K\)。该最后步骤也可取一组含一的 \(K\) 基直接比较张量系数。对任意中间微分域 \(F\)，\(L\) 仍由同一 \(Z\) 在 \(F\) 上生成，且常数不变；\cref{b5:thm:E-pv-existence}说明它是 \(F\) 上的 PV 域。重复这一步得到
\(L^{\operatorname{Gal}_\delta(L/F)}=F\)，而固定 \(F\) 的群在保持代数关系的方程中增加固定各个 \(f\in F\) 的方程，仍为 Zariski 闭群。

另一个方向使用闭子群的有理不变量分离陪集引理：线性代数群 \(G\) 的闭子群 \(H\) 有齐性空间 \(G/H\)，其有理函数能分离一般陪集，故固定全部右 \(H\) 不变有理函数的右平移恰来自 \(H\)。构造这个齐性空间时，取定义 \(H\) 的理想的一个有限维平移稳定生成空间，将其相应子空间或外幂直线嵌入 Grassmann 簇，\(H\) 正是该点的稳定子群。对上述主齐性空间作同样的下降，就得到
\(\operatorname{Gal}_\delta(L/L^H)=H\)。这里省略齐性空间的代数构造和有理函数的下降验证；它们是这一方向的主要技术步骤。

当 \(H\) 正规时，\(G/H\) 是线性代数群，主齐性空间的商可用它的一组忠实矩阵坐标描述。把这些坐标写成解矩阵的有限张量构造，导子给出一个 \(K\) 系数线性系统，其解坐标生成 \(L^H\)，所以该域是 PV 域。反之，若 \(F/K\) 是 PV 扩张，\(G\) 的每个元素都把它的基本解矩阵变成同一矩阵乘一个可逆常数矩阵，故保持 \(F\)。限制自同构因此得到 \(G\to\operatorname{Gal}_\delta(F/K)\)。

取 \(\tau\in\operatorname{Gal}_\delta(F/K)\)，构造扭曲的微分 \(F\) 代数
\[
 {}^\tau L=L\otimes_{F,\tau}F.
\]
这里张量积满足平衡关系
\[
 (f\ell)\otimes b=\ell\otimes\tau(f)b
 \quad(f,b\in F,\ \ell\in L),
\]
其 \(F\) 结构由 \(b\mapsto1\otimes b\) 给出，导子为
\[
 \delta(\ell\otimes b)=\delta(\ell)\otimes b+\ell\otimes\delta(b).
\]
因为 \(\tau\) 是微分自同构，导子与上述平衡关系相容。映射 \(j:L\rightarrow{}^\tau L\)、\(\ell\mapsto\ell\otimes1\) 是微分域同构，并满足 \(j(f\ell)=\tau(f)j(\ell)\)；其逆作为环映射由 \(\ell\otimes b\mapsto\ell\tau^{-1}(b)\) 给出，因而 \({}^\tau L\) 确为域。由于 \(\tau\) 固定 \(K\) 及 \(C\)，矩阵 \(Z^\tau=Z\otimes1\) 仍满足原来的 \(K\) 系数系统，即 \(\delta(Z^\tau)=AZ^\tau\)，它的系数在 \(F\) 上生成 \({}^\tau L\)，且没有新常数。因此 \({}^\tau L/F\) 与 \(L/F\) 都是该系统的 PV 扩张。由 PV 唯一性取得微分 \(F\) 同构 \(\phi:{}^\tau L\xrightarrow{\sim}L\)，则 \(\sigma=\phi\circ j\) 是延拓 \(\tau\) 的微分 \(K\) 自同构。限制映射所以满射，核为 \(H\)，故 \(H\) 正规且商群为所述 Galois 群。这部分省略商主齐性空间的坐标与线性系统的下降。详细的主齐性空间及对应证明见\cite[Section 1.4, Theorem 1.28, Proposition 1.34 and Corollary 1.40]{VanDerPutSinger2003GaloisTheory}，结果的概述见\cite[Theorem 1.3.9 and Proposition 1.3.20, pp. 13--17]{Singer2009DifferentialGalois}。
\end{proofsketch}

对应中的中间域须对导子稳定，子群须为 Zariski 闭子群，两项条件都不可删去。

在\cref{b5:exm:E-logarithm}的 \(K=\mathbb C(x)\)、\(L=K(u)\)、\(u'=1/x\) 中，普通中间域 \(F=K(u^2)\) 就不是微分域。超越性给出 \([L:F]=2\)，而
\[
 (u^2)'=2u/x\notin F.
\]
若这个导数属于 \(F\)，因 \(x\in K\)，就会推出 \(u\in F\)，与次数为二矛盾。另一方面，加法群在特征零的代数闭域上只有零子群和全群两个闭子群：真闭子集有限，而非零元素的整数倍不可能构成有限群。故这个 PV 扩张只有 \(K,L\) 两个中间微分域，普通中间域却远不止这两个。在指数型扩张中则有 \((t^n)'=nt^n\)，所以 \(K(t^n)\) 对导子稳定；同样形式的域生成表达式，须分别检查求导是否仍留在域内。

\ProseParagraphBreak
例如在下一节的乘法群例子中，子群 \(\{2^m:m\in\mathbb Z\}\subset\mathbb C^\times\) 是真子群，却 Zariski 稠密：一元非零多项式不可能有无限多个根。若有理函数 \(f(t)\in K(t)\) 被这些伸缩全部固定，恒等式 \(f(ct)=f(t)\) 清去分母后是关于 \(c\) 的多项式等式，既然在无限多个 \(c\) 上成立，就在所有 \(c\ne0\) 上成立。故该真子群与整个乘法群有相同不动域，不能在删去闭性后仍要求双射。

% !TeX root = ../../Book5.tex
% 纲要标题：### E.3 算例与可解性
% 纲要位置：工作记录/计划纲要.md，第 5086 行。
% 纲要细目（保留原定写作范围）：
%
% * 在 \(\mathbb C(x)\) 上比较指数型、对数型和根式型例子；非齐次方程需要转为齐次系统
% * 用解的代数关系核验乘法群、加法群或有限循环群；对 \(y'=\lambda y/x\) 完整分类并比较单值闭包，不只凭函数名称判断
% * 标准 PV 假设下的 Liouvillian 可解性与群单位连通分支的可解性；以三角系统和代数 Riccati 解解释构造，Kovacic 算法只作为二阶方程专题出口
% #### E.3.1 指数型、对数型与根式型方程
% #### E.3.2 解的代数关系与 Galois 群
% #### E.3.3 Liouvillian 可解性与 Kovacic 算法
%
\section{算例与可解性}
\label{b5:sec:E03}

下面都在 \((K,\delta)=\bigl(\mathbb C(x),\mathrm{d}/\mathrm{d}x\bigr)\) 上讨论。先在有理函数域中用形式变量规定导子，再利用导子的性质证明常数域保持不变；这样构造的扩张可以直接用于后面的微分模计算。

\subsection{指数型、对数型与根式型方程}
\label{b5:subsec:E03:01}

我们需要两个有理函数事实。非零 \(r\in\mathbb C(x)\) 在无穷远处满足 \(r'/r=O(1/x)\)，故其对数导数不可能是非零常数。另一方面，\(r'\) 在有限点的 Laurent 展开没有 \((x-a)^{-1}\) 项：该项只能来自常数项的导数，而它为零。因此 \(r'=1/x\) 没有有理解。

\begin{example}[指数型扩张]\label{b5:exm:E-exponential}
	设 \((K,\delta)=(\mathbb C(x),\mathrm d/\mathrm dx)\)，取在 \(K\) 上超越的元 \(t\)，规定 \(\delta(t)=t\)。则 \(L=K(t)\) 是 \(y'=y\) 的 PV 扩张，微分 Galois 群为 \(\mathbb G_m(\mathbb C)=\mathbb C^\times\)，作用为 \(t\mapsto ct\)。
\end{example}
\begin{proof}
	导子先由 Leibniz 法则延拓到 \(K[t]\)，再按商法则延拓到 \(K(t)\)。每个 \(f\in K(t)\) 在 \(t=0\) 有唯一形式 Laurent 展开
	\[
		f=\sum_{j\ge j_0}a_j(x)t^j\in K(\!(t)\!).
	\]
	这来自分母提出最低次幂后对常数项非零的幂级数求逆，因而展开映射是单射。若 \(f'=0\)，比较系数得 \(a_j'+ja_j=0\)。当 \(j\ne0\) 且 \(a_j\ne0\) 时，这要求 \(a_j'/a_j=-j\)，与无穷远处的计算矛盾；而 \(a_0'=0\) 给出 \(a_0\in\mathbb C\)。故 \(C_L=\mathbb C\)。

	一阶基本解 \(t\) 生成 \(L\)，所以\cref{b5:thm:E-pv-existence}的反向陈述给出 PV 环 \(K[t,t^{-1}]\)。若 \(\sigma\) 是微分自同构，则 \(\bigl(\sigma(t)/t\bigr)'=0\)，故 \(\sigma(t)=ct\)。反过来，\(t\) 的超越性保证每个非零常数伸缩都给出域自同构，且与导子交换。
\end{proof}

\begin{example}[对数型扩张]\label{b5:exm:E-logarithm}
	设 \((K,\delta)=(\mathbb C(x),\mathrm d/\mathrm dx)\)，取在 \(K\) 上超越的元 \(u\)，规定 \(\delta(u)=1/x\)。则 \(L=K(u)\) 是齐次系统
	\[
		Y'=\begin{pmatrix}0&1/x\\0&0\end{pmatrix}Y
	\]
	的 PV 扩张，群为 \(\mathbb G_a(\mathbb C)=(\mathbb C,+)\)，作用为 \(u\mapsto u+c\)。
\end{example}
\begin{proof}
	把非零 \(f\in K(u)\) 在 \(u=\infty\) 展开为 \(f=\sum_{j\le d}a_j(x)u^j\)，其中 \(a_d\ne0\)。若 \(f'=0\)，最高次项给出 \(a_d'=0\)，故 \(a_d\in\mathbb C^\times\)；再比较 \(u^{d-1}\) 项得到
	\[
		a_{d-1}'=-d\,a_d/x.
	\]
	有理函数导数在 \(0\) 的留数为零，故 \(d=0\)。此时 \(a_0\) 为常数；若 \(f-a_0\ne0\)，它仍是常数元，却有严格负的最高次数，重复论证又要求该次数为零，矛盾。因此 \(C_L=\mathbb C\)。

	基本解矩阵可以取
	\[
		Z=\begin{pmatrix}1&u\\0&1\end{pmatrix}.
	\]
	它的行列式为 \(1\)，系数生成 \(L\)，所以由\cref{b5:thm:E-pv-existence}得到 PV 环 \(K[u]\)。自同构满足 \(\bigl(\sigma(u)-u\bigr)'=0\)，必为常数平移；\(u\) 的超越性保证每个这样的平移均可实行。在基本解矩阵上，群元素为 \(\bigl(\begin{smallmatrix}1&c\\0&1\end{smallmatrix}\bigr)\)。
\end{proof}

原式 \(u'=1/x\) 是非齐次标量方程。增加满足 \(v'=0\) 的坐标，并写 \(u'=v/x\)，才得到 PV 定义要求的齐次系统；取 \(v=1\) 恢复原式。

\begin{example}[根式型扩张]\label{b5:exm:E-radical}
	设 \((K,\delta)=(\mathbb C(x),\mathrm d/\mathrm dx)\)，\(n\ge2\)，令 \(t^n=x\)，并将导子延拓到 \(L=K(t)=\mathbb C(t)\)。则 \(L/K\) 是 \(y'=y/(nx)\) 的 PV 扩张，群为 \(\mu_n(\mathbb C)\)，作用为 \(t\mapsto\zeta t\)、\(\zeta^n=1\)。
\end{example}
\begin{proof}
	\(T^n-x\) 对素元 \(x\in\mathbb C[x]\) 满足 Eisenstein 条件\footnote{艾森斯坦（Ferdinand Gotthold Max Eisenstein，1823-1852），德国数学家，研究数论、代数形式和高次互反律。}，所以在 \(K[T]\) 中不可约，\([L:K]=n\)。求导关系得
	\[
		t'=\frac{1}{nt^{n-1}}=\frac{t}{nx}.
	\]
	在 \(\mathbb C(t)\) 中，导子为 \((nt^{n-1})^{-1}\cdot\mathrm{d}/\mathrm{d}t\)，其核为 \(\mathbb C\)，因为特征零有理函数的通常导数为零当且仅当函数为常数。于是 \(t\) 是生成无新常数扩张的基本解，可用\cref{b5:thm:E-pv-existence}。

	自同构必须保持 \(t^n=x\)，故 \(\sigma(t)=\zeta t\)。反之，每个 \(n\) 次单位根均给出固定 \(x\) 的自同构，导子公式表明它与求导交换。因此群恰为所述有限循环群。
\end{proof}

\begin{proposition}\label{b5:prop:E-rank-one-exponent}
	设 \(K=\mathbb C(x)\)、\(\delta=\mathrm d/\mathrm dx\)，\(\lambda\in\mathbb C\)。方程 \(y'=\lambda y/x\) 的 PV 群为：\(\lambda\in\mathbb Z\) 时的平凡群；\(\lambda=p/q\in\mathbb Q\setminus\mathbb Z\)、\(\gcd(p,q)=1,q>1\) 时的 \(\mu_q\)；\(\lambda\notin\mathbb Q\) 时的 \(\mathbb G_m\)。
\end{proposition}
\begin{proof}
	若 \(\lambda=p/q\) 已约分，取 \(t^q=x\)，则 \(y=t^p\) 满足方程。由整数 \(a,b\) 满足 \(ap+bq=1\)，有 \(t=y^ax^b\)，所以 \(K(y)=K(t)\)。当 \(q=1\) 时解已在 \(K\) 中；当 \(q>1\) 时，\cref{b5:exm:E-radical}给出次数为 \(q\) 的无新常数扩张。其全部自同构 \(t\mapsto\zeta t\) 在 \(y\) 上按 \(\zeta^p\) 作用，由互素性这些标量恰为全部 \(q\) 次单位根。

	若 \(\lambda\notin\mathbb Q\)，取超越元 \(t\)，规定 \(t'=\lambda t/x\)。把 \(K(t)\) 中的常数元在 \(t=0\) 处作形式 Laurent 展开 \(\sum_{j\ge j_0}a_j(x)t^j\)，逐项求导并比较系数得
	\[
		a_j'+\frac{j\lambda}{x}a_j=0.
	\]
	非零有理函数在 \(x=0\) 可以写成 \(a_j=x^m b\)，其中 \(m\in\mathbb Z\)，\(b\) 在零点正则且非零。故 \(a_j'/a_j=m/x+b'/b\) 的留数为整数。对 \(j\ne0\)，等式要求它等于 \(-j\lambda\notin\mathbb Z\)，不可能；只有 \(a_0\in\mathbb C\) 留下。因此无新常数，\cref{b5:thm:E-pv-existence}给出 PV 扩张。全部非零常数伸缩都与导子相容，所以群为 \(\mathbb G_m\)。
\end{proof}

在有理指数情形，若错误地保留一个超越解变量 \(y\)，则 \((y^q/x^p)'=0\) 会产生新常数。PV 条件须将它固定为某个 \(c\in\mathbb C^\times\)，于是得到 \(y^q=cx^p\)；因常数域代数闭，伸缩解可将 \(c\) 化为一。解析地，局部解 \(x^\lambda\) 绕零点一周的乘子为 \(e^{2\pi i\lambda}\)。有理指数给出相应有限循环群，非有理指数给出无限阶元素；后者的循环子群在 \(\mathbb G_m\) 中 Zariski 稠密。因此在这个秩一正则奇点算例中，PV 群就是单值群的 Zariski 闭包，这个直接计算不替代一般正则奇点方程的对应定理。

\subsection{解的代数关系与 Galois 群}
\label{b5:subsec:E03:02}

基本解矩阵的常数变换必须保持其代数关系。在例~\ref{b5:exm:E-radical}中，\(t^n=x\) 把伸缩限制为 \(c^n=1\)；在例~\ref{b5:exm:E-exponential}中，没有这种关系，全部非零伸缩便都存在。

\ProseParagraphBreak
对系统 \(Y'=\operatorname{diag}(1,2)Y\)，沿用 \(t'=t\)，可取
\[
	Z=\operatorname{diag}(t,t^2).
\]
第二个非零对角元是第一个的平方，所以群不是两个独立乘法群，而是
\[
	\bigl\{\operatorname{diag}(c,c^2):c\in\mathbb C^\times\bigr\}.
\]
这里也能核验全部关系：从 \(K[T_1^{\pm1},T_2^{\pm1}]\) 代入 \((t,t^2)\) 的核恰为 \((T_2-T_1^2)\)，因为商环就是 \(K[T_1,T_1^{-1}]\)，而 \(t\) 超越。

\begin{proposition}\label{b5:prop:E-exponential-fixed-field}
	设 \((K,\delta)=(\mathbb C(x),\mathrm d/\mathrm dx)\)，\(t\) 在 \(K\) 上超越且 \(\delta(t)=t\)，\(L=K(t)\)。其微分 Galois 群 \(\mathbb C^\times\) 以 \(t\mapsto ct\) 作用。对每个 \(n\ge1\)，子群 \(\mu_n=\{c\in\mathbb C^\times:c^n=1\}\) 的不动域为 \(K(t^n)\)，整个 \(\mathbb C^\times\) 的不动域为 \(K\)。
\end{proposition}
\begin{proof}
	\(1,t,\ldots,t^{n-1}\) 是 \(K(t^n)\) 上的基：以超越元 \(s=t^n\) 对 \(T^n-s\) 用 Eisenstein 判别即可。取本原 \(n\) 次单位根 \(\zeta\)，若 \(\sum_{j=0}^{n-1}b_jt^j\) 被 \(t\mapsto\zeta t\) 固定，比较基系数便得 \(b_j=0\) 对 \(j>0\) 成立。

	若 \(f\) 被全部伸缩固定，在 \(t=0\) 的 Laurent 展开中便有 \(a_jc^j=a_j\) 对所有 \(c\ne0\) 成立。对 \(j\ne0\) 可取 \(c^j\ne1\)，故只留下常数项，\(f\in K\)。
\end{proof}

由于 \((t^n)'=nt^n\)，\(K(t^n)\) 对导子稳定，确实是\cref{b5:thm:E-galois-correspondence}所要求的中间微分域。

\subsection{Liouvillian 可解性与 Kovacic 算法}
\label{b5:subsec:E03:03}

\begin{definition}\label{b5:def:E-liouvillian}
	设 \((K,\delta)\) 为特征零微分域。给定有限微分域塔
	\[
		K=K_0\subset K_1\subset\cdots\subset K_r,\qquad K_{i+1}=K_i(a_i)\quad(0\le i<r),
	\]
	若每步满足下列条件之一：\(a_i\) 在 \(K_i\) 上代数、\(\delta(a_i)\in K_i\)、或 \(a_i\ne0\) 且 \(\delta(a_i)/a_i\in K_i\)，则称为 \term[Liouvillian-ta]{Liouvillian 塔}[Liouvillian tower]。\footnote{刘维尔（Joseph Liouville，1809-1882），法国数学家，研究分析与数论，构造了早期的超越数例子，并研究微分方程的积分表达。} 在常数域 \(C=\ker\delta\) 代数闭的情形，对 \(n\ge1\)、\(A\in M_n(K)\)，设 \(L/K\) 为系统 \(\delta(Y)=AY\) 的 PV 扩张。若存在 \(L\) 到这样的塔顶的微分 \(K\) 嵌入，则称该系统 \term[Liouvillian-kejie]{Liouvillian 可解}[Liouvillian solvable]。
\end{definition}

后两类步骤分别对应积分和积分的指数。第一类允许一般代数扩张，并不要求由根式组成；这解释了下面为何只出现单位连通分支。

\begin{theorem}[Liouvillian 可解性准则]\label{b5:thm:E-liouvillian-criterion}
	设 \((K,\delta)\) 为特征零微分域，常数域 \(C=\ker\delta\) 代数闭，\(n\ge1\)，\(A\in M_n(K)\)，\(L/K\) 为系统 \(\delta(Y)=AY\) 的 PV 扩张，\(G=\operatorname{Gal}_\delta(L/K)\)。该系统 Liouvillian 可解，当且仅当 \(C\) 上线性代数群 \(G\) 的单位连通分支 \(G^\circ\) 可解。
\end{theorem}

\begin{proofsketch}
	先设 \(G^\circ\) 可解。由\cref{b5:thm:E-galois-correspondence}，\(F=L^{G^\circ}\) 是 \(K\) 上的 PV 扩张，其群为有限群 \(G/G^\circ\)，这里使用 PV 扩张的超越次数等于其代数群维数的定理，见\cite[Proposition 1.34(3)]{VanDerPutSinger2003GaloisTheory}，所以 \(\operatorname{trdeg}_K F=\dim(G/G^\circ)=0\)。PV 域由有限多个元素生成，所以 \(F/K\) 是有限代数扩张。Lie--Kolchin 定理说明，\footnote{科尔钦（Ellis Robert Kolchin，1916-1991），美国数学家，研究微分代数与微分域的 Galois 理论。}代数闭域上的连通可解线性代数群在有限维有理表示中保持一条完全旗；将它用于 \(G^\circ\) 在解空间上的表示。这里使用的是群的三角化定理，见\cite[Theorem 1.5.8, p. 44]{Singer2009DifferentialGalois}。

	在原 PV 域 \(L\) 中选取各列适配这条旗的基本解矩阵 \(Z\)。前 \(j\) 列所张成的 \(L\) 子空间在 \(G^\circ\) 的作用下稳定。将这些列转置后取行最简形矩阵 \(H\)。对 \(\sigma\in G^\circ\)，逐项施加 \(\sigma\) 固定 \(0,1\)，因而 \(\sigma(H)\) 仍为行最简形；子空间的稳定性又使它与 \(H\) 具有相同的行空间。由行最简形的唯一性，\(\sigma(H)=H\)，所以每个系数被 \(G^\circ\) 固定，从而落在 \(L^{G^\circ}=F\)。这些列又是水平向量，对其线性组合施加 \(\nabla=\delta-A\) 时，只对组合系数求导，故各子空间对联络稳定。联络及下降后的基均定义在 \(F\) 上，于是得到 \(F\) 上的微分子模旗。选取适配旗的 \(F\) 基，令换基矩阵为 \(B\in\operatorname{GL}_n(F)\)，就使
	\[
		T=B^{-1}Z,\qquad \widetilde A=B^{-1}AB-B^{-1}\cdot\delta(B),\qquad \delta(T)=\widetilde A T
	\]
	中的 \(T\) 与 \(\widetilde A\) 都为上三角矩阵。

	二阶情形说明上三角结构为何允许逐层求解。若系统矩阵为
	\(\bigl(\begin{smallmatrix}a&b\\0&d\end{smallmatrix}\bigr)\)，先取 \(t_1'=at_1,t_2'=dt_2\)。基本解矩阵可写为
	\[
		T=\begin{pmatrix}t_1&z\\0&t_2\end{pmatrix},
		\qquad z'=az+bt_2, \qquad (z/t_1)'=b\,t_2/t_1.
	\]
	所以取得两个对角解后，只需再加入右侧的一个原函数。例如 \(a=1,d=2,b=1/x\) 时，\(z=e^x\int e^x/x\,\mathrm dx\)；最后的原函数不必初等，仍是 Liouvillian 塔允许的步骤。

	现在完全在 \(L\) 内构造塔。先加入各个非零对角元 \(t_i=T_{ii}\)，它们满足 \(\delta(t_i)/t_i=\widetilde A_{ii}\in F\)，是定义中允许的指数步骤。随后逐列处理 \(T_{ij}\)，在第 \(j\) 列中按 \(i=j-1,j-2,\ldots,1\) 的次序进行。变常数公式给出
	\[
		\delta(T_{ij}/t_i) =\frac1{t_i}\sum_{k=i+1}^j\widetilde A_{ik}T_{kj}.
	\]
	右侧只涉及此前已经加入的元素，所以加入 \(T_{ij}/t_i\) 是一个原函数步骤。这样取得 \(T\) 的全部条目，而 \(L=F(T_{ij})\)，塔顶就是原来的 \(L\)。每层均为 \(L\) 的微分子域，常数域均为 \(C\)。

	反向设 \(L\) 嵌入塔顶 \(E\)。先处理塔可能引入的新常数。令 \(C_1=\ker(\delta:E\to E)\)，在 \(E\) 的代数闭包 \(\overline E\) 中取 \(C_1\) 的代数闭包 \(C_2\)，并唯一延拓导子。若 \(z\in\overline E\) 满足 \(\delta(z)=0\)，将它在 \(E\) 上的首一最小多项式求导，所得多项式次数更低且仍以 \(z\) 为根，故所有系数的导数为零。于是 \(z\) 在 \(C_1\) 上代数，属于 \(C_2\)。反过来，\(C_2\) 中元素的导数均为零，所以 \(\overline E\) 的常数恰为 \(C_2\)。每个复合域 \(K_iC_2\) 因而有同一代数闭常数域 \(C_2\)，原来的生成元仍给出一个有限 Liouvillian 塔。

	同一基本解矩阵在 \(KC_2\) 上生成 \(LC_2\)，且后者无新常数，由\cref{b5:thm:E-pv-existence}它是 PV 扩张。这里的群正是 \(G\) 的常数域扩张：若 \(a_1,\ldots,a_r\in C_2\) 在 \(C\) 上线性无关，却有项数最少的非零 \(L\) 线性关系，除以第一个非零系数可写成 \(a_1+\sum_{i=2}^r\ell_i a_i=0\)。求导得 \(\sum_{i=2}^r\delta(\ell_i)a_i=0\)，第一项已经消失；若任一导数非零，就与项数最少矛盾。因此全部 \(\ell_i\) 都在 \(C_L=C\) 中，原关系又与 \(C\) 线性无关矛盾。故 \(L\) 与 \(C_2\) 在 \(C\) 上线性无交。记原 PV 环为 \(R\)，先扩张常数，再将 \(K\otimes_C C_2\) 的非零元素取逆，得到嵌入 \(LC_2\) 的环
	\[
		R_2=(R\otimes_C C_2)\otimes_{K\otimes_C C_2}KC_2 =KC_2[Z_{ij},(\det Z)^{-1}].
	\]
	反向 PV 判据说明它是 \(LC_2/KC_2\) 的 PV 环。对\cref{b5:thm:E-galois-correspondence}证明中的主齐性空间恒等式同时扩张常数并作上述取逆，得到
	\[
		R_2\otimes_{KC_2}R_2 \simeq R_2\otimes_{C_2}C_2[G_{C_2}],
	\]
	因而微分 Galois 群为 \(G_{C_2}\)。代数闭常数域扩张保持单位分支，其可解性也可由导出列经忠实平坦基变换检验。因此只须证明塔没有新常数的情形。

	以下在上述常数基变换后对塔长归纳，将共同常数域、基域与 PV 域仍记为 \(C\)、\(K\)、\(L\)。塔长为零时 \(L=K\)。记首步为 \(K(t)\)，令
	\[
		E_0=L\cap K(t),\qquad H=\operatorname{Gal}_\delta\bigl(L(t)/K(t)\bigr).
	\]
	\(L(t)/K(t)\) 是同一系统的 PV 扩张，而余下的塔长少一，故归纳假设给出 \(H^\circ\) 可解。限制到 \(L\) 将 \(H\) 识别为 \(G\) 的闭子群；因为
	\[
		L^H=L(t)^H\cap L=K(t)\cap L=E_0,
	\]
	Galois 对应又给出 \(H=\operatorname{Gal}_\delta(L/E_0)\)。若 \(E_0=K\)，则 \(H=G\)，结论已得。若 \(t\) 在 \(K\) 上代数，\(E_0/K\) 有限，\(H\) 是有限指数闭子群，故 \(H^\circ=G^\circ\)。

	剩下 \(t\) 超越的两种情形。若 \(\delta(t)\in K\)，则 \(K(t)/K\) 是加法群型的 PV 扩张。加法群在特征零域上没有非平凡真闭子群，所以非平凡交域只能是 \(E_0=K(t)\)。若 \(\delta(t)/t\in K\)，则 \(K(t)/K\) 是乘法群型的 PV 扩张；由乘法群的闭子群分类，非平凡交域为 \(E_0=K(t^d)\)，其中 \(d\ge1\)。前一情形中的 \(t\) 在 \(G\) 作用下只相差常数，后一情形中的 \(t^d\) 只相差非零常数倍，故 \(E_0\) 对 \(G\) 稳定。Galois 对应给出正规子群 \(H\)，并给出
	\[
		G/H\cong\mathbb G_a\quad\text{或}\quad G/H\cong\mathbb G_m.
	\]
	商群交换，所以 \(G^\circ\) 的导出子群包含于 \(H\)；连通代数群的导出子群连通，因而实际上包含于 \(H^\circ\)。由 \(H^\circ\) 可解，便得 \(G^\circ\) 可解；这一群论步骤也见\cite[Corollary A.38]{VanDerPutSinger2003GaloisTheory}。

	共同常数域下的交域归纳见\cite[Section 1.5, Theorem 1.43]{VanDerPutSinger2003GaloisTheory}。准则与塔的定义亦见\cite[Definition 1.5.6 and Theorem 1.5.7, pp. 43--44]{Singer2009DifferentialGalois}。
\end{proofsketch}

乘法群、加法群均可解，有限群的单位连通分支平凡，所以三个算例都符合准则。即使有限微分 Galois 群本身不可解，相应代数解仍属于允许的代数步骤；把条件误写成“\(G\) 可解”会排除这些情形。另外，“存在一个非零 Liouvillian 解”弱于整个系统可解，本定理说的是足够组成基本解矩阵的全部解。

\ProseParagraphBreak
对二阶方程 \(y''=r(x)y\)、\(r\in\mathbb C(x)\)，非零解的对数导数 \(v=y'/y\) 满足 Riccati 方程
\[
	v'+v^2=r.
\]
若找到有理 \(v\)，则 \(y=\exp\int v\,dx\) 是 Liouvillian 解，另一解可通过降阶得到。Kovacic 算法还处理适当代数扩域中的候选对数导数，以判定此类二阶方程的 Liouvillian 可解性；实行算法要求常数域中的必要代数运算可以有效进行。其有限分类与次数界见 \cite[Section 1.5, pp. 49--51]{Singer2009DifferentialGalois}。这里没有给出完整算法，尤其不能因找不到有理 \(v\) 就断言不可解。

\ProseParagraphBreak
例如 \(r=x+5/(16x^2)\) 时，\(v=\sqrt x-1/(4x)\) 满足
\[
	v'=\frac1{2\sqrt x}+\frac1{4x^2},\qquad v^2=x-\frac1{2\sqrt x}+\frac1{16x^2},\qquad v'+v^2=r.
\]
它给出 Liouvillian 解 \(y_1=x^{-1/4}\exp(\tfrac23x^{3/2})\)，另一解为 \(y_2=y_1\int y_1^{-2}\,\mathrm dx\)。后者满足方程，且 Wronskian 为 \(y_1^2(y_2/y_1)'=1\)，所以整个系统有 Liouvillian 基本解。可是该 Riccati 方程没有有理解：若 \(v\in\mathbb C(x)\) 的多项式部分次数为 \(d\ge1\)，则 \(v^2\) 的最高次数为 \(2d\)，\(v'\) 的次数至多为 \(d-1\)，不能得到次数为一的 \(r\)；若 \(v\) 在无穷远有界，则 \(v'+v^2\) 也有界，同样不可能。真正的对数导数位于代数扩域中，寻找有理解的失败不足以判定 Liouvillian 不可解。

% !TeX root = ../../Book5.tex
% 纲要标题：### E.4 张量结构与解空间上的表示
% 纲要位置：工作记录/计划纲要.md，第 5098 行。
% 纲要细目（保留原定写作范围）：
%
% * 微分模的张量不变量与 Tannaka 重构思想；由秩一张量幂恢复有限群关系，非分裂对数模型比较水平截面的正合性与表示的半单性；参数化理论转第三卷附录 E.4
% * 比较形式局部基域对 PV 群的影响；非正则秩一方程区分单值与指数因子。行列式准则及二阶环面模型说明群约束与群的实际确定；Stokes 与 Riemann–Hilbert 方法转附录 D 及分析／几何教材
% * Picard–Vessiot 存在唯一性在 E.2 证明；可计算 Galois 群和更高阶算法可继续阅读专著
% #### E.4.1 张量不变量与 Tannaka 重构
% #### E.4.2 形式局部基域与解析延拓
% #### E.4.3 行列式方程与群的限制
% ---
%
\section{张量结构与解空间上的表示}
\label{b5:sec:E04}

微分 Galois 群在原方程解空间上的作用，可以延伸到张量积、交替积和商空间。方程已有的配对和张量关系也会约束群。以下先讨论张量不变量，再比较改变基域对微分 Galois 群的影响。

\subsection{张量不变量与 Tannaka 重构}
\label{b5:subsec:E04:01}

设 \(K\) 的特征为零，常数域 \(C\) 代数闭，\(M\) 为有限维微分模，\(L/K\) 为其系统的 PV 域。记 \(\langle M\rangle_\otimes\) 为由 \(M\) 和平凡微分模 \(K\)，经有限直和、张量积、对偶、子模及商模得到的微分模所成的全子范畴。对其中的对象 \(N\)，取
\[
 \omega(N)=\ker\bigl(\nabla:L\otimes_KN\longrightarrow L\otimes_KN\bigr).
\]
它是 \(C\) 向量空间。微分自同构 \(\sigma\) 只作用于 \(L\) 因子，因与导子交换而保持水平截面；于是 \(\omega(N)\) 得到 \(G=\operatorname{Gal}_\delta(L/K)\) 的表示。张量积上的作用和通常的表示张量积相符。

\ProseParagraphBreak
上一节的系统 \(Z=\operatorname{diag}(t,t^2)\)、\(\delta t=t\)，已经给出一个非线性的张量关系。相应的两个秩一微分模分别记为 \(M_1=Ke_1\)、\(M_2=Ke_2\)，规定 \(\nabla e_1=-e_1\)、\(\nabla e_2=-2e_2\)，则 \(te_1,t^2e_2\) 为水平基，且映射
\[
\phi:M_1\otimes M_1\longrightarrow M_2,\qquad
e_1\otimes e_1\longmapsto e_2
\]
与联络相容。若纤维函子的一个张量自同构分别在 \(\omega(M_1),\omega(M_2)\) 上作用为标量 \(c_1,c_2\)，张量相容性使其在 \(\omega(M_1)^{\otimes2}\) 上作用为 \(c_1^2\)，对 \(\phi\) 的自然性便迫使 \(c_2=c_1^2\)。所以前面的矩阵群方程 \(\operatorname{diag}(c,c^2)\)，也可由微分模之间的张量同构恢复。

\ProseParagraphBreak
有限循环群也可以由张量同构读出。对\cref{b5:exm:E-radical}的 \(K=\mathbb C(x)\)、\(t^n=x\)，取秩一微分模 \(M=Ke\)，规定 \(\nabla e=-e/(nx)\)。其水平基为 \(te\)，而
\[
 \nabla(e^{\otimes n})=-\frac1x e^{\otimes n},
 \qquad \nabla(xe^{\otimes n})=0.
\]
因此 \(M^{\otimes n}\simeq K\)，同构 \(\phi\) 将 \(e^{\otimes n}\) 送到 \(x^{-1}\)，于是将水平向量 \((te)^{\otimes n}\) 送到一。若张量自同构在 \(\omega(M)\) 上作用为标量 \(c\)，对 \(\phi\) 的自然性及单位对象上的恒等作用迫使 \(c^n=1\)。这先给出 \(G\subseteq\mu_n\)；根式型算例又证明每个单位根都由 \(t\mapsto\zeta t\) 实现，故为相等。张量平方关系联系两个伸缩参数，张量幂平凡化则迫使一个参数有限阶。

这里保留的资料不只是一个交换范畴：\(\omega\) 要与其中全部态射自然相容，并保持张量积、对偶与单位对象。第21章中，普通代数同构不能恢复群基元，而 Hopf 结构能借 \(\Delta(g)=g\otimes g\) 恢复它们；这里的张量积与单位对象保留的就是同类资料。若只保留对象间的加性态射而遗忘这些张量关系，便没有理由从它重构原群。

\begin{theorem}[PV 理论的张量表述]\label{b5:thm:E-tannaka}
设 \((K,\delta)\) 为特征零微分域，其常数域 \(C=\ker\delta\) 代数闭；\((M,\nabla_M)\) 为 \(K\) 上的有限维微分模，\(L/K\) 为其线性系统的 PV 扩张，\(G=\operatorname{Gal}_\delta(L/K)\)。令 \(\langle M\rangle_\otimes\) 为有限维微分 \(K\) 模范畴中，由 \(M\) 及平凡微分模 \((K,\delta)\) 经有限直和、张量积、对偶及子商生成的全子范畴。对其中的对象 \((N,\nabla_N)\)，在 \(L\otimes_KN\) 上取延拓联络 \(\nabla(\ell\otimes v)=\delta(\ell)\otimes v+\ell\otimes\nabla_N(v)\)，并令
\(\omega(N)=\ker\nabla\)。则 \(\omega\) 将 \(\langle M\rangle_\otimes\) 等价地映到 \(C\) 上线性代数群 \(G\) 的有限维有理表示范畴。它保持张量积，并且是忠实、正合的；每个 \(N\) 都满足
\(\dim_C\omega(N)=\dim_KN\)。对任意交换 \(C\) 代数 \(S\)，有自然同构 \(G(S)\simeq\operatorname{Aut}^\otimes(S\otimes_C\omega)\)，因而 \(G\) 可由这一范畴连同函子 \(\omega\) 恢复。
\end{theorem}
\symidx{\(\langle M\rangle_\otimes,\omega\)：微分模生成的张量子范畴与水平截面纤维函子}

\begin{proof}
\(L\otimes_KM\) 由基本解矩阵给出水平基，因而是平凡微分模。直和、张量与对偶显然保持这一性质，子商也保持：在水平基下，设 \(W\subset L^r\) 是对逐系数求导稳定的子空间，取它唯一的行最简形基。每个基向量 \(v_i\) 求导后主元位置全为零；由求导稳定性，写 \(v_i'=\sum_j a_jv_j\)。在 \(v_j\) 的主元位置，右端坐标恰为 \(a_j\)，故全部 \(a_j=0\)，即 \(v_i'=0\)，所以所有非主元系数均在 \(C\) 中。于是 \(W\) 有水平基，商也有水平基。由此，对所有 \(N\in\langle M\rangle_\otimes\)，自然映射
\[
 \eta_N:L\otimes_C\omega(N)\longrightarrow L\otimes_KN,
 \qquad \ell\otimes s\longmapsto\ell s
\]
是同构，给出维数等式及张量相容性。微分模同态在水平基下是常数矩阵；对短正合列取这些矩阵，秩在 \(C\) 与 \(L\) 上相同，所以取水平截面保持满射，也保持正合性。张量 \(L\) 是忠实的，故 \(\omega\) 忠实。

设 \(T:\omega(N_1)\to\omega(N_2)\) 为 \(G\) 等变的 \(C\) 线性映射，即 \(T(gv)=gT(v)\)。把它用 \(\eta_{N_1},\eta_{N_2}\) 延伸为 \(L\) 线性映射。这个映射在原来的 \(K\) 基下的矩阵系数被 \(G\) 固定，由\cref{b5:thm:E-galois-correspondence}属于 \(K\)，故下降成 \(K\) 线性映射 \(f:N_1\to N_2\)。其标量扩张由下面的交换图唯一确定：
\[
\begin{tikzcd}[column sep=large,row sep=large]
L\otimes_C\omega(N_1)\arrow[r,"1\otimes T"]
\arrow[d,"\eta_{N_1}","\sim"']
&L\otimes_C\omega(N_2)\arrow[d,"\eta_{N_2}","\sim"']\\
L\otimes_KN_1\arrow[r,dashed,"1\otimes f"']&L\otimes_KN_2
\end{tikzcd}
\]
在水平基下，\(1\otimes T\) 的系数在 \(C\) 中，因而与逐系数求导相容；两个 \(\eta_{N_i}\) 也与延拓联络相容。所以 \(1\otimes f\) 与两端的联络相容。差 \(f\nabla_{N_1}-\nabla_{N_2}f\) 是 \(K\) 线性映射，张量到 \(L\) 后为零，由忠实平坦性可得
\[
 f\nabla_{N_1}=\nabla_{N_2}f.
\]
故 \(f\) 是微分模同态。反向取水平截面即可，证明 \(\omega\) 全忠实。

还要说明全部有理表示都能得到。记 \(V=\omega(M)\)，它是 \(G\) 的忠实表示。\(C[G]\) 由 \(V\) 的矩阵坐标及逆行列式生成，因而每个有限维、平移稳定的坐标函数空间都是有限个 \(V,V^*\) 的张量幂的子商。任意有限维有理表示 \(U\) 用各坐标泛函的矩阵系数嵌入有限份 \(C[G]\)：在群单位元处求值便说明这个嵌入单射，有限个系数又落在一个有限维坐标函数空间中。因此 \(U\) 是这些张量幂直和的子商。

设相应的 \(G\) 稳定子空间为 \(W\subseteq\omega(N)\)，通过上面的 \(\eta_N\)，把 \(L\otimes_C W\) 识别为 \(L\otimes_KN\) 的子空间。它由水平向量张成，故对联络稳定。\(G\) 在后者上的半线性作用为 \(\sigma(\ell\otimes n)=\sigma(\ell)\otimes n\)；\(\eta_N\) 与此作用相容，且 \(W\) 为 \(G\) 稳定，故这个 \(L\) 子空间也为 \(G\) 稳定。在原 \(K\) 基下作行最简化，唯一性使每个矩阵系数被 \(G\) 固定，故均属于 \(K\)。这个子空间于是下降为微分子模，商同样下降。这样每个 \(U\) 都由某个 \(N\) 得到，证明范畴等价。

最后固定任意交换 \(C\) 代数 \(S\)。\(S\otimes_C\omega\) 的张量自同构 \(\alpha\) 由各对象上的 \(\alpha_N\) 组成，它们是与态射、张量积及单位对象相容的 \(S\) 线性可逆映射。记其在 \(V=\omega(M)\) 上的分量为
\[
 g=\alpha_M\in\operatorname{GL}(S\otimes_CV).
\]
每个对象都是 \(V,V^*\) 的有限张量构造的子商，而 \(S\) 在 \(C\) 上平坦；自然性与张量相容性因此使 \(g\) 唯一决定全部 \(\alpha_N\)。

选定 \(V\) 的一组基，令 \(n=\dim_CV\)，并写
\[
 A=C[\operatorname{GL}(V)]=C[X_{ij},\det(X)^{-1}],
 \qquad I=\ker\bigl(A\longrightarrow C[G]\bigr).
\]
取 \(I\) 的有限生成元 \(f_1,\ldots,f_r\)。清到共同的行列式分母后，可以选取整数 \(d,N\ge0\)，使这些生成元都属于有限维空间
\[
 E=\det(X)^{-N}C[X_{ij}]_{\le d}.
\]
\(E\) 对右平移 \((r_hf)(x)=f(xh)\) 稳定，因为
\[
 r_h(X_{ij})=\sum_kX_{ik}h_{kj},
 \qquad r_h\bigl(\det(X)^{-N}\bigr)=\det(h)^{-N}\det(X)^{-N}.
\]
矩阵坐标的每一行在右平移下是 \(V\) 的一份拷贝，而 \(\det(X)^{-N}\) 对应 \((\bigwedge^nV^*)^{\otimes N}\)，所以 \(E\) 是 \(V,V^*\) 的有限张量构造的子商。于是 \(\alpha\) 在 \(S\otimes_CE\) 上的作用正是右平移 \(r_g\)。

由于 \(G\) 是闭子群，右平移的余作用 \(A\to A\otimes_CC[G]\) 将 \(I\) 送入 \(I\otimes_CC[G]\)。它也保持 \(E\)，故 \(W=I\cap E\) 是 \(G\) 子表示。由前面证明的范畴等价，它对应一个微分子模；\(\alpha\) 的自然性遂给出
\[
 r_g(S\otimes_CW)=S\otimes_CW.
\]
在单位矩阵处的求值映射 \(\varepsilon:A\to C\) 零化 \(I\)，因为单位矩阵属于 \(G\)。将它扩张为 \(\varepsilon_S:S\otimes_CA\to S\)，对每个 \(f_i\in W\) 有
\[
 f_i(g)=\varepsilon_S\bigl(r_g(1\otimes f_i)\bigr)=0.
\]
因此由矩阵 \(g\) 定义的坐标环映射 \(A\to S\) 零化 \(I\)，因子化通过 \(C[G]\)，这就是 \(g\in G(S)\)。反过来，每个 \(g\in G(S)\) 在有理表示上的作用都给出张量自同构。上述构造与 \(C\) 代数同态 \(S\to S'\) 相容，得到所述群函子的自然同构，即\term[Tannaka-chonggou]{Tannaka 重构}[Tannaka reconstruction]。\footnote{淡中忠郎（假名：たんなか ただお，罗马音：Tadao Tannaka，1908-1986），日本数学家，研究表示论与对偶理论；以其名字命名的重构思想从张量表示恢复群。}

PV 范畴等价的标准表述见\cite[Theorem 1.5.2, p. 41]{Singer2009DifferentialGalois}及\cite[Theorem 2.33]{VanDerPutSinger2003GaloisTheory}；任意系数代数上的群函子重构见\cite[Theorem B.20]{VanDerPutSinger2003GaloisTheory}。
\end{proof}

正合性依赖于这些微分模被同一 PV 域平凡化；在一般微分域中，取水平截面只有左正合性，不能仅由“取一个算子的核”推出上述结论。

\begin{proposition}\label{b5:prop:E-logarithmic-extension}
在 \((K,\delta)=(\mathbb C(x),\mathrm d/\mathrm dx)\) 上，设 \(M=Ke_1\oplus Ke_2\)，\(\nabla e_1=0,\nabla e_2=-e_1/x\)。则
\(0\to Ke_1\to M\to K\to0\) 是平凡微分模的非分裂扩张，且 \(H(M)=\mathbb Ce_1\)，所以取水平截面不右正合。若 \(L=K(u)\)、\(u'=1/x\) 为对数型 PV 域，则 \(\omega(M)\) 有水平基
\(v_1=e_1,v_2=ue_1+e_2\)，群 \(\mathbb G_a\) 在这组基上按
\(\bigl(\begin{smallmatrix}1&c\\0&1\end{smallmatrix}\bigr)\) 作用，所诱导的短正合列仍在群表示范畴中不分裂。
\end{proposition}
\begin{proof}
商中的 \(\overline e_2\) 水平。微分模的分裂须将一送到 \(e_2+fe_1\)，要求 \(f'=1/x\)；有理函数导数在零点的留数为零，故这样的 \(f\in K\) 不存在。水平向量 \(ae_1+be_2\) 则须满足 \(b'=0,a'=b/x\)，于是 \(b=0,a\in\mathbb C\)。其像在商的水平空间中为零，未能保持满射。

扩域后，直接求导验证 \(v_1,v_2\) 水平，其坐标矩阵行列式为一，所以它们是基。自同构 \(u\mapsto u+c\) 将 \(v_2\) 送到 \(v_2+cv_1\)。因而向量空间上的序列
\(0\to\mathbb C\to\omega(M)\to\mathbb C\to0\) 正合；但是商中一的任何提升 \(v_2+av_1\) 都不被全部平移固定，故不存在群等变的分裂。
\end{proof}

\(L\) 中的水平基使扩张成为平凡微分模，却把原来的扩张信息保留在群作用中。纤维函子正合只保证底层向量空间的短正合列仍正合，其线性截面未必与群作用相容，因而不能推出表示范畴半单。第22.7节推出分裂时，使用的是 \((-)^G\) 正合，从而使 \(\operatorname{Hom}_G(U,-)\simeq(U^*\otimes-)^G\) 保持满射，将 \(1_U\) 提升为等变截面。这里的正合纤维函子没有这项取不变量的条件；对数例中的平移作用正好阻止这样的提升。

\ProseParagraphBreak
对具体张量，群的限制常可直接核验。例如系统 \(Y'=AY\) 上若存在可逆矩阵 \(B\in M_n(K)\)，满足
\begin{equation}\label{b5:eq:E-horizontal-pairing}
 B'+A^{\mathsf T}B+BA=0,
\end{equation}
则对基本解矩阵 \(Z\) 求导给出 \((Z^{\mathsf T}BZ)'=0\)。故
\(Q=Z^{\mathsf T}BZ\in M_n(C)\) 可逆。任意 \(\sigma(Z)=ZC_\sigma\) 都满足
\[
 C_\sigma^{\mathsf T}QC_\sigma=Q,
\]
因为 \(\sigma\) 固定 \(B\) 和 \(Q\)。当 \(B\) 对称或交替时，这分别把群限制在相应正交群或辛群内。这只证明包含关系；若还有别的张量关系，实际群可能更小。\cref{b5:thm:E-tannaka}说明为什么必须同时考虑所有这些构造。

\subsection{形式局部基域与解析延拓}
\label{b5:subsec:E04:02}

一个方程的群依赖于所选基域。例~\ref{b5:exm:E-exponential}中 \(y'=y\) 在 \(\mathbb C(x)\) 上的群是 \(\mathbb G_m\)，但把基域扩大为 \(\mathbb C(\!(x)\!)\) 后，
\[
 t_0=\sum_{m\ge0}\frac{x^m}{m!}
\]
已是该域中的可逆基本解。形式求导给出 \(t_0'=t_0\)，而 Laurent 级数的导数为零恰好意味着它是常数。因此这个局部形式基域上的 PV 扩张就是基域自身，群平凡。全局的代数关系和一个普通点附近的形式关系不能混为一谈。

\ProseParagraphBreak
对 \(u'=1/x\) 则不同：\(\mathbb C(\!(x)\!)\) 中任何元素的导数仍没有 \(x^{-1}\) 项，故不能在基域内找到 \(u\)。取形式超越元 \(u\) 并规定 \(u'=1/x\)，例~\ref{b5:exm:E-logarithm}中的证明仍成立：其中只用了基域常数为 \(\mathbb C\) 及“导数留数为零”，这两个性质对 Laurent 级数同样成立。因此局部形式群仍是加法群。若改在穿孔圆盘的通用覆盖上解释解析解，一次绕原点的延拓又把 \(\log x\) 变成 \(\log x+2\pi i\)，给出群中的一个具体元素。

\ProseParagraphBreak
还须区别形式群与解析单值。在 \(K=\mathbb C(\!(x)\!)\) 上考虑 \(y'=y/x^2\)。非零 \(f\in K\) 可写成 \(x^m b(x)\)，其中 \(b\in\mathbb C[\![x]\!]^\times\)，所以
\[
 \frac{f'}f=\frac mx+\frac{b'}b,
 \qquad \frac{b'}b\in\mathbb C[\![x]\!].
\]
对数导数至多有一阶极点，不可能等于 \(1/x^2\)。取超越元 \(t\) 满足 \(t'=t/x^2\)，在 \(K(t)\) 的 \(t\)-Laurent 展开中，常数元的非零系数须满足 \(a_j'/a_j=-j/x^2\)；上述极点限制排除所有 \(j\ne0\)，所以仍无新常数。它是 PV 扩张，全部非零常数伸缩都存在，形式局部群为 \(\mathbb G_m\)。解析解 \(e^{-1/x}\) 却在穿孔圆盘上单值，解析单值群平凡。它的本性奇点不能用负幂有下界的形式 Laurent 级数表示；这里非平凡的形式群来自指数因子，并不需要非平凡 Stokes 跳跃。

\ProseParagraphBreak
这些例子区分了有理系数、形式局部和解析延拓三种问题。一般非正则奇点还需要扇形渐近解和 Stokes 数据，具体说明见\appref{b5:app:D}。只有形式级数而没有解析收敛或渐近条件时，尚不能谈论这种延拓。

\subsection{行列式方程与群的限制}
\label{b5:subsec:E04:03}

\begin{proposition}[行列式方程]\label{b5:prop:E-wronskian}
设 \((K,\delta)\) 为特征零微分域，常数域 \(C=\ker\delta\) 代数闭，\(n\ge1\)，\(A\in M_n(K)\)，\(L/K\) 为系统 \(\delta(Y)=AY\) 的 PV 扩张，\(Z\in\operatorname{GL}_n(L)\) 为基本解矩阵，撇号表示 \(\delta\)。则
\[
 (\det Z)'=\operatorname{tr}(A)\det Z.
\]
由 \(\sigma(Z)=ZC_\sigma\) 定义 \(G=\operatorname{Gal}_\delta(L/K)\) 的群矩阵。则 \(G\subseteq\operatorname{SL}_n(C)\) 当且仅当 \(\det Z\in K^\times\)，亦当且仅当 \(\operatorname{tr}(A)=f'/f\) 对某个 \(f\in K^\times\) 成立；后一条件下 \(\det Z=cf\)、\(c\in C^\times\)。
\end{proposition}
\begin{proof}
按行列式的多线性逐列求导，得到
\((\det Z)'=\operatorname{tr}(Z^{-1}Z')\det Z\)。代入 \(Z'=AZ\) 并利用矩阵迹在共轭下不变，便得第一式。若 \(\operatorname{tr}A=f'/f\)，则商 \(\det Z/f\) 的导数为零，无新常数给出 \(\det Z=cf\)，故 \(\det Z\in K^\times\)。此时自同构固定 \(\det Z\)，对 \(\sigma(Z)=ZC_\sigma\) 取行列式便得 \(\det C_\sigma=1\)。反之，若所有 \(C_\sigma\) 的行列式为一，\(\det Z\) 就被全群固定；\cref{b5:thm:E-galois-correspondence}给出 \(L^G=K\)，故 \(\det Z\in K^\times\)，取 \(f=\det Z\) 即得迹的条件。
\end{proof}

特别地，\(y''=r(x)y\) 的一阶系统矩阵为
\(\bigl(\begin{smallmatrix}0&1\\r&0\end{smallmatrix}\bigr)\)，迹为零，所以微分 Galois 群包含在 \(\operatorname{SL}_2(C)\) 中。这给出上一节二阶算法所研究的群范围。由最高外幂的联络公式，这个判别也等价于 \(\bigwedge^nM\) 是平凡微分模：Wronskian 正是最高外幂水平方程的一个非零解。

\begin{proposition}\label{b5:prop:E-second-order-torus}
在 \((K,\delta)=(\mathbb C(x),\mathrm d/\mathrm dx)\) 上，取超越元 \(t\) 满足 \(t'=t\)。方程 \(y''=y\) 的 PV 域为 \(K(t)\)，基本解矩阵可取为
\[
 Z=\begin{pmatrix}t&t^{-1}\\t&-t^{-1}\end{pmatrix},
 \qquad \det Z=-2.
\]
其微分 Galois 群矩阵恰为
\(\{\operatorname{diag}(c,c^{-1}):c\in\mathbb C^\times\}\)，是 \(\operatorname{SL}_2(\mathbb C)\) 的真子群。
\end{proposition}
\begin{proof}
矩阵的两列分别由解 \(t,t^{-1}\) 及其导数组成，行列式非零，条目生成 \(K(t)\)；\cref{b5:exm:E-exponential}证明该域无新常数。
系统矩阵 \(A=\bigl(\begin{smallmatrix}0&1\\1&0\end{smallmatrix}\bigr)\) 的两条常数特征直线给出两个微分子模。取常数矩阵
\(P=\bigl(\begin{smallmatrix}1&1\\1&-1\end{smallmatrix}\bigr)\)，有 \(P^{-1}AP=\operatorname{diag}(1,-1)\)；这两个子模在纤维空间中也给出两条分别稳定的直线，所以群矩阵必须对角。由 \(\det Z\in K\)，\cref{b5:prop:E-wronskian}进一步要求两个对角元的乘积为一。最后，每个 \(t\mapsto ct\) 都是微分自同构，作用于两列的标量分别为 \(c,c^{-1}\)，于是上述上界恰为实际群。
\end{proof}

这个计算依次使用了行列式和微分子模两个约束，再验证每个允许的伸缩都存在。只检查行列式不能把 \(\operatorname{SL}_2\) 缩小到这个一维环面。

\ProseParagraphBreak
配对和行列式是可实际计算的张量不变量；它们给出群的上界，而找全解之间的代数关系才能确定群。多参数理论还要加入与主导子交换的其他导子，所得群一般由微分多项式方程刻画，已不属于本附录的普通线性代数群情形。相关概念放在第三卷附录 E.4；有效群计算和高阶算法可继续阅读 \cite[Sections 1.3--1.5]{Singer2009DifferentialGalois}。


\begin{chaptersummary}
微分模中的 Leibniz 法则使换基公式多出导数项，也决定了张量和对偶的作用。PV 理论在代数闭常数域上选取由基本解矩阵生成且无新常数的扩张，把微分域的中间结构与代数群的闭子群相联系。具体群必须由全部代数关系核验，不能仅凭解的函数名称判断。

\ProseParagraphBreak
Liouvillian 塔允许代数扩张、积分与积分的指数，所以判别条件落在群的单位连通分支上。张量重构从表示论解释这种群如何编码微分模；实际计算时，配对和行列式常先给出可核验的限制。全局有理函数域、形式局部域和解析函数域的选择会改变问题，使用这些结果前应始终写清基域和常数域。
\end{chaptersummary}

\BookExercises

\begin{exercise}
设 \((K,\delta)\) 是微分域。由 Leibniz 法则推出 \(\delta(1)=0\)、逆元公式及商法则。再说明特征 \(p>0\) 时 \(\bigl(\mathbb F_p(x),\mathrm{d}/\mathrm{d}x\bigr)\) 的常数域不是 \(\mathbb F_p\)。
\end{exercise}
\begin{solution}
对 \(1\cdot1\) 求导得 \(\delta(1)=2\delta(1)\)，故 \(\delta(1)=0\)。对 \(aa^{-1}=1\) 求导得
\(\delta(a^{-1})=-a^{-2}\cdot\delta(a)\)；再对 \(ab^{-1}\) 求导即可得
\(\delta(a/b)=\bigl(\delta(a)b-a\cdot\delta(b)\bigr)/b^2\)。
在正特征中，\((x^p)'=0\)，所以常数域含 \(\mathbb F_p(x^p)\)，严格大于 \(\mathbb F_p\)。本附录的特征零假设排除了这一现象。
\end{solution}

\begin{exercise}
在微分域 \((K,\delta)=(\mathbb C(x),\mathrm d/\mathrm dx)\) 上考虑 \(y'=y/x\)。以换基矩阵 \(P=x\) 计算新系统，并求原微分模的水平基。
\end{exercise}
\begin{solution}
由 \eqref{b5:eq:E-gauge-change}，
\(\widetilde A=P^{-1}AP-P^{-1}P'=1/x-1/x=0\)。
若原基 \(e\) 满足 \(\nabla(e)=-e/x\)，则
\(\nabla(xe)=e+x(-e/x)=0\)，所以 \(xe\) 是水平基。这个方程已经在基域中平凡化，PV 群为平凡群。
\end{solution}

\begin{exercise}
设 \((K,\delta)\) 为特征零微分域，微分模 \(M,N\) 在基下的水平系统分别为 \(v'=Av,w'=Bw\)。把 \(\operatorname{Hom}_K(M,N)=M^*\otimes_KN\) 的元素写成矩阵 \(T\)，证明它水平当且仅当 \(T'=BT-TA\)。在两个模均为秩一时，用此条件判定非零微分模同态何时存在，并解释取 \(N=(K,\delta)\) 如何恢复对偶方程。
\end{exercise}
\begin{solution}
对映射应用对偶与张量联络公式，有 \((\nabla T)(m)=\nabla_N(Tm)-T\nabla_M(m)\)。在基坐标中它的矩阵为 \(T'-BT+TA\)，所以水平条件正是所述方程。
若 \(A=a,B=b\) 为标量，则非零同态由某个 \(g\in K^\times\) 给出，条件为 \(g'/g=b-a\)。这也是两个秩一模同构的条件，因为非零的一维线性映射可逆。取 \(N=(K,\delta)\)，即 \(B=0\)，行矩阵满足 \(T'=-TA\)，转置后即得 \(w'=-A^{\mathsf T}w\)。
\end{solution}

\begin{exercise}
设 \(K=\mathbb C(x)\)、\(\delta=\mathrm d/\mathrm dx\)，\(g\in K\)，微分模 \(M_g=Ke_1\oplus Ke_2\) 满足 \(\nabla e_1=0,\nabla e_2=-ge_1\)。证明它是平凡微分模的分裂扩张，当且仅当 \(g\) 在每个有限点的留数都为零。分别计算分裂与非分裂情形的 \(\dim_{\mathbb C}H(M_g)\)。
\end{exercise}
\begin{solution}
分裂须以水平向量 \(e_2+fe_1\) 提升商中的一，等价于 \(f'=g\) 有有理解。有理函数导数在每个有限点的留数都为零。反之，对 \(g\) 作部分分式分解；若全部简单极点项的系数为零，则多项式项和 \((x-a)^{-j}\)、\(j\ge2\) 各有有理原函数，所以 \(g=f'\)。
水平条件为 \(b'=0,a'=bg\)。若存在上述 \(f\)，则 \(e_1,e_2+fe_1\) 是水平基，维数为二；若不存在，非零常数 \(b\) 会给出 \((a/b)'=g\)，故只能有 \(b=0,a\in\mathbb C\)，维数为一。
\end{solution}

\begin{exercise}
在 \(K=\mathbb C(x)\) 上令 \(M=Ke_1\oplus Ke_2\oplus Ke_3\)，联络满足 \(\nabla e_1=0,\nabla e_2=-e_1/x,\nabla e_3=-2e_1/x\)。取对数型 PV 域 \(L=K(u)\)、\(u'=1/x\)，求一组水平基及群作用矩阵。再在基域上分解 \(M\)，说明两个非零耦合项为何没有给出两个独立的平移参数。
\end{exercise}
\begin{solution}
可取 \(v_1=e_1,v_2=ue_1+e_2,v_3=2ue_1+e_3\)。直接求导可知三者水平，坐标矩阵为单位上三角矩阵。\(u\mapsto u+c\) 的作用矩阵为
\[
 \begin{pmatrix}1&c&2c\\0&1&0\\0&0&1\end{pmatrix}.
\]
在基域中 \(e_3-2e_2\) 已水平，故
\(M=(Ke_1\oplus Ke_2)\oplus K(e_3-2e_2)\) 是非分裂的二维对数模与一个平凡模的直和。两个耦合系数的比例已在基域中固定，它们使用同一个原函数 \(u\)，群只含一个平移参数。
\end{solution}

\begin{exercise}
在 \(K=\mathbb C(x)\) 上取超越元 \(t\) 满足 \(t'=t\)，令 \(L=K(t)\)。求方程 \(y''=4y\) 的 PV 子域 \(F\subseteq L\)、基本解矩阵及群矩阵。求限制映射 \(\operatorname{Gal}_\delta(L/K)\to\operatorname{Gal}_\delta(F/K)\) 的核，并比较两个群都同构于 \(\mathbb G_m\) 与 \(F\ne L\) 是否矛盾。
\end{exercise}
\begin{solution}
令 \(s=t^2\)，则 \(s'=2s\)，可取
\[
 Z=\begin{pmatrix}s&s^{-1}\\2s&-2s^{-1}\end{pmatrix},
 \qquad \det Z=-4.
\]
条目生成 \(F=K(s)\)，它作为 \(L\) 的子域无新常数，所以是 PV 域。自同构 \(s\mapsto ds\) 对应 \(\operatorname{diag}(d,d^{-1})\)，每个 \(d\in\mathbb C^\times\) 均可取。限制映射把 \(t\mapsto ct\) 送到 \(s\mapsto c^2s\)，核为 \(\mu_2\)，且 \([L:F]=2\)。两个抽象群同构，不意味着这里的限制满射也是单射；本例的非平凡核已直接说明这一点，也不能由抽象群同构推出 PV 域相等。
\end{solution}

\begin{exercise}
设 \((K,\delta)=(\mathbb C(x),\mathrm d/\mathrm dx)\)，\(n\ge2\)，\(t^n=x\)，将导子延拓到 \(L=K(t)\)。取正整数 \(d\) 并令 \(g=\gcd(n,d)\)。求 \(K(t^d)\) 在根式 PV 扩张 \(L/K\) 中的对应子群及其对基域的次数。
\end{exercise}
\begin{solution}
固定 \(t^d\) 的自同构须满足 \(\zeta^n=\zeta^d=1\)，故子群为 \(\mu_g\)，阶为 \(g\)。取整数 \(a,b\) 使 \(ad+bn=g\)，则
\(t^g=(t^d)^ax^b\)，而 \(t^d\) 又是 \(t^g\) 的幂，所以中间域为 \(K(t^g)\)。
其生成元满足 \((t^g)^{n/g}=x\)，Eisenstein 判别给出次数 \(n/g\)，与群的指数一致。
\end{solution}

\begin{exercise}
设 \((K,\delta)=(\mathbb C(x),\mathrm d/\mathrm dx)\)，\(u\) 在 \(K\) 上超越且 \(\delta(u)=1/x\)，扩张 \(K(u)/K\) 无新常数。若 \(v\in K(u)\) 满足 \(v'=2/x\)，证明 \(v=2u+c\)，其中 \(c\in\mathbb C\)。用此结论求系统矩阵为 \(\bigl(\begin{smallmatrix}0&1/x\\0&0\end{smallmatrix}\bigr)\) 与 \(\bigl(\begin{smallmatrix}0&2/x\\0&0\end{smallmatrix}\bigr)\) 的块对角和时的实际微分 Galois 群。
\end{exercise}
\begin{solution}
\((v-2u)'=0\)，而常数域为 \(\mathbb C\)，故 \(v=2u+c\)。选择基本解时可令 \(v=2u\)。两个块上的平移参数必须分别为 \(a,2a\)，不能独立选择；群为一个加法群，在两块中分别以
\(\bigl(\begin{smallmatrix}1&a\\0&1\end{smallmatrix}\bigr)\) 和
\(\bigl(\begin{smallmatrix}1&2a\\0&1\end{smallmatrix}\bigr)\) 作用。
\end{solution}

\begin{exercise}
取相互独立的超越元 \(t,c\)，在 \(L=\mathbb C(x)(t,c)\) 上规定 \(t'=t,c'=0\)。说明为什么这个含有基本解 \(t\) 的域不是 \(y'=y\) 的 PV 扩张。
\end{exercise}
\begin{solution}
常数域至少含 \(\mathbb C(c)\)，已超过基域常数 \(\mathbb C\)，违反无新常数条件。另外基本解 \(t\) 本身只生成 \(\mathbb C(x)(t)\)，无法生成另一个独立元 \(c\)。存在某个基本解矩阵不是定义的全部内容。
\end{solution}

\begin{exercise}
设 \((K,\delta)=(\mathbb C(x),\mathrm d/\mathrm dx)\)，\(u\) 在 \(K\) 上超越且 \(\delta(u)=1/x\)。在对数型 PV 扩张 \(K(u)/K\) 中，证明整数平移子群
\(\{u\mapsto u+m:m\in\mathbb Z\}\) 的不动域也为 \(K\)。这是否违反 Galois 对应？
\end{exercise}
\begin{solution}
若 \(f(u+1)=f(u)\)，在 \(u=\infty\) 展开。若非零最高次数为 \(d\ne0\)，则 \(f(u+1)-f(u)\) 的最高非零项为 \(d a_du^{d-1}\)，矛盾。若最高次数为零，先减去常数项；只要还有非零项，其最高次数严格负，再用同一论证便矛盾。因此 \(f\in K\)。
整数子群是加法群的 Zariski 稠密真子群，不是闭子群，故不在\cref{b5:thm:E-galois-correspondence}要求双射的对象中。
\end{solution}

\begin{exercise}
对 \(y''=(x^2+1)y\)，找一个 Liouvillian 解，并用它写出第二个解的积分表达式。证明两解线性无关。
\end{exercise}
\begin{solution}
Riccati 方程 \(v'+v^2=x^2+1\) 有解 \(v=x\)，故可取 \(y_1=e^{x^2/2}\)。令
\[
 y_2=y_1\int y_1^{-2}\,dx
     =e^{x^2/2}\int e^{-x^2}\,dx.
\]
写 \(y_2=y_1w\)，代入原方程后剩下
\(y_1w''+2y_1'w'=0\)，由 \(w'=y_1^{-2}\) 可核验。Wronskian
\(y_1y_2'-y_1'y_2=y_1^2w'=1\)，故两解无关。所用操作属于 Liouvillian 塔，尽管积分未必是初等函数。
\end{solution}

\begin{exercise}
设 \((K,\delta)\) 为特征零微分域，常数域 \(C=\ker\delta\) 代数闭，某线性系统的 PV 扩张 \(L/K\) 的微分 Galois 群为有限不可解群。按本附录定义，该系统是否 Liouvillian 可解？它的解是否因此必能用根式表达？
\end{exercise}
\begin{solution}
有限代数群的单位连通分支平凡，故由\cref{b5:thm:E-liouvillian-criterion}，系统 Liouvillian 可解。定义允许任意代数步骤，而不是仅允许根式步骤，所以不能由此推出根式表达；有限不可解 Galois 群恰好说明这两种可解性条件不同。
\end{solution}

\begin{exercise}
在常数域为代数闭域 \(C\) 的特征零微分域 \(K\) 上，设
\[
 A=\begin{pmatrix}0&a\\-a&0\end{pmatrix},\qquad a\in K.
\]
证明相应 PV 群保持某个非退化对称型且行列式为 \(1\)。能否仅由此断定群等于完整的特殊正交群？
\end{exercise}
\begin{solution}
取 \(B=I\)，有 \(B'=0\)、\(A^{\mathsf T}+A=0\)，所以
\eqref{b5:eq:E-horizontal-pairing}成立。于是 \(Q=Z^{\mathsf T}Z\) 为非退化常数对称型，群保持 \(Q\)。同时 \(\operatorname{tr}A=0\)，\cref{b5:prop:E-wronskian}给出行列式为 \(1\)。
只能得到包含关系：若 \(a=0\)，系统已在基域中有常数基本解，群平凡，而二维非退化型的特殊正交群并不平凡。
\end{solution}

\begin{exercise}
在 \(K=\mathbb C(x)\) 上取整数 \(n\ge2\)，考虑系统
\[
 Y'=\operatorname{diag}(1/x,1/(nx))Y.
\]
求其 PV 域与微分 Galois 群，并求行列式特征标的像。能否用 \(K\) 上的规范变换使系统矩阵的迹为零？
\end{exercise}
\begin{solution}
取 \(t^n=x\)，则基本解矩阵为 \(Z=\operatorname{diag}(x,t)\)。它生成根式型 PV 域 \(K(t)\)，群矩阵为 \(\operatorname{diag}(1,\zeta)\)、\(\zeta^n=1\)。行列式特征标的像是整个 \(\mu_n\)，因而群不包含于 \(\operatorname{SL}_2\)。
原矩阵的迹为 \((1+1/n)/x\)。若规范变换能使迹为零，由取迹公式就须有 \((\det P)'/\det P=(1+1/n)/x\)。写 \(\det P=x^m g\)，其中 \(m\in\mathbb Z\)，且 \(g\) 在 \(0\) 处正则并非零，则 \((\det P)'/\det P=m/x+g'/g\)。后项在 \(0\) 处正则，故留数为整数 \(m\)；但 \(1+1/n\) 不是整数，故这样的 \(P\in\operatorname{GL}_2(K)\) 不存在。这也与\cref{b5:prop:E-wronskian}的判别一致。
\end{solution}

\begin{exercise}
比较 \(y'=y\) 在 \(\mathbb C(x)\) 与 \(\mathbb C(\!(x)\!)\) 上的 PV 群，并解释为什么这不与基本解矩阵换基只引起共轭相矛盾。
\end{exercise}
\begin{solution}
前者由例~\ref{b5:exm:E-exponential}得乘法群；后者已经含有可逆基本解
\(\sum_{m\ge0}x^m/m!\)，且 Laurent 级数域常数为 \(\mathbb C\)，故 PV 扩张平凡，群也平凡。\cref{b5:prop:E-group-matrix}比较的是同一个基域和同一个 PV 扩张中的基本解矩阵，此处则扩大了基域，假设不同。
\end{solution}

\begin{optionalexercise}
设 \(\lambda\in\mathbb C\setminus\mathbb Q\)，取在 \(K=\mathbb C(x)\) 上代数独立的 \(s,t\)，规定 \(s'=s,t'=\lambda t\)。证明 \(K(s,t)\) 无新常数，并求 \(Y'=\operatorname{diag}(1,\lambda)Y\) 的 PV 群。
\end{optionalexercise}
\begin{solution}
将有理函数域嵌入迭代 Laurent 级数域 \(K(\!(s)\!)(\!(t)\!)\)。展开中的系数 \(a_{ij}(x)\) 若对应常数元，便满足
\[
 a_{ij}'+(i+\lambda j)a_{ij}=0.
\]
有理函数的非零常数对数导数不存在，所以非零系数只能出现在
\(i+\lambda j=0\)。由于 \(i,j\in\mathbb Z\) 且 \(\lambda\notin\mathbb Q\)，这迫使 \(i=j=0\)，相应系数属于 \(\mathbb C\)。故无新常数。
基本解矩阵 \(\operatorname{diag}(s,t)\) 生成该域。微分自同构 \(\sigma\) 保持 \(s'=s\)、\(t'=\lambda t\)，故 \((\sigma s/s)'=(\sigma t/t)'=0\)；由已证的常数域结论，只能分别作伸缩 \(\sigma s=cs\)、\(\sigma t=dt\)，其中 \(c,d\in\mathbb C^\times\)。反过来，代数独立性使任意 \((s,t)\mapsto(cs,dt)\) 唯一延拓为 \(K(s,t)\) 的域自同构，逆映射由 \(c^{-1},d^{-1}\) 给出。它在 \(K,s,t\) 上与导子交换，因而在整个域上交换，群因此为 \(\mathbb G_m^2\)。
\end{solution}
